\begin{thebibliography}{999}
\small
\addcontentsline{toc}{chapter}{Bibliography}
\markboth{Bibliography}{Bibliography}

\bibitem{abrahamdeo2023} K. Abraham and N. Deo. Deep Gaussian process priors for Bayesian inference in nonlinear inverse problems. Preprint, arXiv:2312.14294, 2023.
\bibitem{aeschbacher2012} S. Aeschbacher, M. A. Beaumont and A. Futschik. A novel approach for choosing summary statistics in approximate Bayesian computation. \emph{Genetics} 192 (2012), 1027--1047.
\bibitem{albergo2023} M. S. Albergo and E. Vanden Eijnden. Building normalizing flows with stochastic interpolants. \emph{International Conference on Learning Representations} (2023).
\bibitem{albertrath2020} C. G. Albert and K. Rath. Gaussian process regression for data fulfilling linear differential equations with localized sources. \emph{Entropy} 22 (2020), 152.
\bibitem{abc2022} D. Albritton, E. Bru\'e and M. Colombo. Nonuniqueness of Leray solutions of the forced Navier Stokes equations. \emph{Annals of Mathematics} 196 (2022), 415--455.
\bibitem{alexander1975} R. Alexander. The infinite hard sphere system. M.S. thesis, University of California, Berkeley, 1975.
\bibitem{alili1997} L. Alili, D. Dufresne and M. Yor. Sur l'identit\'e de Bougerol pour les fonctionnelles exponentielles du mouvement brownien avec drift. In \emph{Exponential functionals and principal values related to Brownian motion}, ed. M. Yor, pages 3--14. Biblioteca de la Revista Matem\'atica Iberoamericana, 1997.
\bibitem{alon1998} N. Alon, M. Krivelevich and B. Sudakov. Finding a large hidden clique in a random graph. \emph{Random Structures and Algorithms} 13 (1998), 457--466.
\bibitem{alpoge2026} L. Alp\"oge, T. Buckmaster and M. P. Coiculescu. Extending the C\'ordoba Mart\'inez Zoroa IPM blowup to uniformly spacetime smooth forcing. Preprint, arXiv:2609.16470, 2026.
\bibitem{altmeyer2022} R. Altmeyer. Polynomial time guarantees for sampling based posterior inference in high dimensional generalised linear models. Preprint, arXiv:2208.13296, 2022.
\bibitem{altschuler2017} J. Altschuler, J. Weed and P. Rigollet. Near linear time approximation algorithms for optimal transport via Sinkhorn iteration. \emph{Advances in Neural Information Processing Systems} 30 (2017), 1961--1971.
\bibitem{anderson1982} B. D. O. Anderson. Reverse time diffusion equation models. \emph{Stochastic Processes and their Applications} 12 (1982), 313--326.
\bibitem{andrewsmallows1974} D. F. Andrews and C. L. Mallows. Scale mixtures of normal distributions. \emph{Journal of the Royal Statistical Society B} 36 (1974), 99--102.
\bibitem{adh2010} C. Andrieu, A. Doucet and R. Holenstein. Particle Markov chain Monte Carlo methods. \emph{Journal of the Royal Statistical Society B} 72 (2010), 269--342.
\bibitem{akp1990} D. Applegate, R. Kannan and N. Polson. Random polynomial time algorithms for sampling from joint distributions. Technical report, Carnegie Mellon University, 1990.
\bibitem{applegate1991} D. Applegate and R. Kannan. Sampling and integration of near log concave functions. In \emph{Proceedings of the Twenty Third Annual ACM Symposium on Theory of Computing}, pages 156--163. ACM Press, 1991.
\bibitem{arimoto1972} S. Arimoto. An algorithm for computing the capacity of arbitrary discrete memoryless channels. \emph{IEEE Transactions on Information Theory} 18 (1972), 14--20.
\bibitem{arnold1957} V. I. Arnold. On functions of three variables. \emph{Doklady Akademii Nauk SSSR} 114 (1957), 679--681; English translation in \emph{American Mathematical Society Translations}, Series 2, 28 (1963), 51--54.
\bibitem{arnold1963} V. I. Arnold. Proof of a theorem of A. N. Kolmogorov on the invariance of quasiperiodic motions under small perturbations of the Hamiltonian. \emph{Russian Mathematical Surveys} 18 (1963), 9--36.
\bibitem{arnoldavez1968} V. I. Arnold and A. Avez. \emph{Ergodic Problems of Classical Mechanics}. W. A. Benjamin, 1968.
\bibitem{asano1970} T. Asano. Theorems on the partition functions of the Heisenberg ferromagnets. \emph{Journal of the Physical Society of Japan} 29 (1970), 350--359.
\bibitem{aspect1982} A. Aspect, J. Dalibard and G. Roger. Experimental test of Bell's inequalities using time varying analyzers. \emph{Physical Review Letters} 49 (1982), 1804--1807.
\bibitem{bachelier1900} L. Bachelier. Th\'eorie de la sp\'eculation. \emph{Annales scientifiques de l'\'Ecole Normale Sup\'erieure} 17 (1900), 21--86.
\bibitem{bahadurzabell1979} R. R. Bahadur and S. L. Zabell. Large deviations of the sample mean in general vector spaces. \emph{The Annals of Probability} 7 (1979), 587--621.
\bibitem{bain2009} A. Bain and D. Crisan. \emph{Fundamentals of Stochastic Filtering}. Springer, 2009.
\bibitem{bakry1985} D. Bakry and M. \'Emery. Diffusions hypercontractives. In \emph{S\'eminaire de Probabilit\'es XIX}, pages 177--206. Lecture Notes in Mathematics 1123, Springer, 1985.
\bibitem{bakry2014} D. Bakry, I. Gentil and M. Ledoux. \emph{Analysis and geometry of Markov diffusion operators}. Grundlehren der mathematischen Wissenschaften 348, Springer, 2014.
\bibitem{bandeira2023} A. S. Bandeira, A. Maillard, R. Nickl and S. Wang. On free energy barriers in Gaussian priors and failure of cold start MCMC for high dimensional unimodal distributions. \emph{Philosophical Transactions of the Royal Society A} 381 (2023), 20220150.
\bibitem{barak2019} B. Barak, S. Hopkins, J. Kelner, P. K. Kothari, A. Moitra and A. Potechin. A nearly tight sum of squares lower bound for the planted clique problem. \emph{SIAM Journal on Computing} 48 (2019), 687--735.
\bibitem{baraud2016} Y. Baraud. Bounding the expectation of the supremum of an empirical process over a (weak) VC major class. \emph{Electronic Journal of Statistics} 10 (2016), 1709--1728.
\bibitem{barbier2019} J. Barbier and N. Macris. The adaptive interpolation method: a simple scheme to prove replica formulas in Bayesian inference. \emph{Probability Theory and Related Fields} 174 (2019), 1133--1185.
\bibitem{barbier2019b} J. Barbier, F. Krzakala, N. Macris, L. Miolane and L. Zdeborov\'a. Optimal errors and phase transitions in high dimensional generalized linear models. \emph{Proceedings of the National Academy of Sciences} 116 (2019), 5451--5460.
\bibitem{barbier2019c} J. Barbier and N. Macris. The adaptive interpolation method for proving replica formulas. Applications to the Curie Weiss and Wigner spike models. \emph{Journal of Physics A: Mathematical and Theoretical} 52 (2019), 294002.
\bibitem{bardostartar1973} C. Bardos and L. Tartar. Sur l'unicit\'e r\'etrograde des \'equations paraboliques et quelques questions voisines. \emph{Archive for Rational Mechanics and Analysis} 50 (1973), 10--25.
\bibitem{bgl1991} C. Bardos, F. Golse and D. Levermore. Fluid dynamic limits of kinetic equations. I. Formal derivations. \emph{Journal of Statistical Physics} 63 (1991), 323--344.
\bibitem{bgl1993} C. Bardos, F. Golse and C. D. Levermore. Fluid dynamic limits of kinetic equations II: convergence proofs for the Boltzmann equation. \emph{Communications on Pure and Applied Mathematics} 46 (1993), 667--753.
\bibitem{bardostiti2018} C. Bardos and E. S. Titi. Onsager's conjecture for the incompressible Euler equations in bounded domains. \emph{Archive for Rational Mechanics and Analysis} 228 (2018), 197--207.
\bibitem{barnes2012} C. P. Barnes, S. Filippi, M. P. H. Stumpf and T. Thorne. Considerate approaches to constructing summary statistics for ABC model selection. \emph{Statistics and Computing} 22 (2012), 1181--1197.
\bibitem{barron1986} A. R. Barron. Entropy and the central limit theorem. \emph{The Annals of Probability} 14 (1986), 336--342.
\bibitem{barrowgreen1997} J. Barrow Green. \emph{Poincar\'e and the three body problem}. American Mathematical Society, 1997.
\bibitem{batchelor1969} G. K. Batchelor. Computation of the energy spectrum in homogeneous two dimensional turbulence. \emph{Physics of Fluids} 12 (1969), II 233--II 239.
\bibitem{bayati2011} M. Bayati and A. Montanari. The dynamics of message passing on dense graphs, with applications to compressed sensing. \emph{IEEE Transactions on Information Theory} 57 (2011), 764--785.
\bibitem{bayerdiaconis1992} D. Bayer and P. Diaconis. Trailing the dovetail shuffle to its lair. \emph{The Annals of Applied Probability} 2 (1992), 294--313.
\bibitem{bkm1984} J. T. Beale, T. Kato and A. Majda. Remarks on the breakdown of smooth solutions for the 3D Euler equations. \emph{Communications in Mathematical Physics} 94 (1984), 61--66.
\bibitem{beaumont2002} M. A. Beaumont, W. Zhang and D. J. Balding. Approximate Bayesian computation in population genetics. \emph{Genetics} 162 (2002), 2025--2035.
\bibitem{beckner1975} W. Beckner. Inequalities in Fourier analysis. \emph{Annals of Mathematics} 102 (1975), 159--182.
\bibitem{bcpw2019} J. Bedrossian, M. Coti Zelati, S. Punshon Smith and F. Weber. A sufficient condition for the Kolmogorov 4/5 law for stationary martingale solutions to the 3D Navier Stokes equations. \emph{Communications in Mathematical Physics} 367 (2019), 1045--1075.
\bibitem{bell1964} J. S. Bell. On the Einstein Podolsky Rosen paradox. \emph{Physics Physique Fizika} 1 (1964), 195--200.
\bibitem{benaim1999} M. Bena\"im. Dynamics of stochastic approximation algorithms. In \emph{S\'eminaire de Probabilit\'es XXXIII}, pages 1--68. Lecture Notes in Mathematics, Springer, 1999.
\bibitem{benarous2020} G. Ben Arous, R. Gheissari and A. Jagannath. Algorithmic thresholds for tensor PCA. \emph{The Annals of Probability} 48 (2020), 2052--2087.
\bibitem{benes1981} V. E. Bene\v{s}. Exact finite dimensional filters for certain diffusions with nonlinear drift. \emph{Stochastics} 5 (1981), 65--92.
\bibitem{bennett1982} C. H. Bennett. The thermodynamics of computation, a review. \emph{International Journal of Theoretical Physics} 21 (1982), 905--940.
\bibitem{berard2014} J. B\'erard, P. Del Moral and A. Doucet. A lognormal central limit theorem for particle approximations of normalizing constants. \emph{Electronic Journal of Probability} 19 (2014), 94.
\bibitem{bernardo1979} J. M. Bernardo. Reference posterior distributions for Bayesian inference. \emph{Journal of the Royal Statistical Society B} 41 (1979), 113--147.
\bibitem{bernton2019} E. Bernton, P. E. Jacob, M. Gerber and C. P. Robert. Approximate Bayesian computation with the Wasserstein distance. \emph{Journal of the Royal Statistical Society B} 81 (2019), 235--269.
\bibitem{berti2004} P. Berti, L. Pratelli and P. Rigo. Limit theorems for a class of identically distributed random variables. \emph{The Annals of Probability} 32 (2004), 2029--2052.
\bibitem{bertini2015} L. Bertini, A. De Sole, D. Gabrielli, G. Jona Lasinio and C. Landim. Macroscopic fluctuation theory. \emph{Reviews of Modern Physics} 87 (2015), 593--636.
\bibitem{berut2012} A. B\'erut, A. Arakelyan, A. Petrosyan, S. Ciliberto, R. Dillenschneider and E. Lutz. Experimental verification of Landauer's principle linking information and thermodynamics. \emph{Nature} 483 (2012), 187--189.
\bibitem{besag1995} J. Besag, P. Green, D. Higdon and K. Mengersen. Bayesian computation and stochastic systems (with discussion). \emph{Statistical Science} 10 (1995), 3--66.
\bibitem{beurling1960} A. Beurling. An automorphism of product measures. \emph{Annals of Mathematics} 72 (1960), 189--200.
\bibitem{bhatnagar1954} P. L. Bhatnagar, E. P. Gross and M. Krook. A model for collision processes in gases. I. Small amplitude processes in charged and neutral one component systems. \emph{Physical Review} 94 (1954), 511--525.
\bibitem{bbm1975} I. Bia\l{}ynicki Birula and J. Mycielski. Uncertainty relations for information entropy in wave mechanics. \emph{Communications in Mathematical Physics} 44 (1975), 129--132.
\bibitem{bpy2001} P. Biane, J. Pitman and M. Yor. Probability laws related to the Jacobi theta and Riemann zeta functions, and Brownian excursions. \emph{Bulletin of the American Mathematical Society} 38 (2001), 435--465.
\bibitem{biane2009} P. Biane. Matrix valued Brownian motion and a paper by P\'olya. In \emph{S\'eminaire de Probabilit\'es XLII}, pages 171--185. Lecture Notes in Mathematics 1979, Springer, 2009.
\bibitem{billingsley1999} P. Billingsley. \emph{Convergence of probability measures}. Second edition, Wiley, 1999.
\bibitem{birgepolson2012} J. R. Birge and N. G. Polson. Optimisation via slice sampling. Preprint, arXiv:1212.2135, 2012.
\bibitem{birkhoff1931} G. D. Birkhoff. Proof of the ergodic theorem. \emph{Proceedings of the National Academy of Sciences} 17 (1931), 656--660.
\bibitem{birkhoff1957} G. Birkhoff. Extensions of Jentzsch's theorem. \emph{Transactions of the American Mathematical Society} 85 (1957), 219--227.
\bibitem{blachman1965} N. M. Blachman. The convolution inequality for entropy powers. \emph{IEEE Transactions on Information Theory} 11 (1965), 267--271.
\bibitem{blackwell1973} D. Blackwell and J. B. MacQueen. Ferguson distributions via P\'olya urn schemes. \emph{The Annals of Statistics} 1 (1973), 353--355.
\bibitem{blahut1972} R. E. Blahut. Computation of channel capacity and rate distortion functions. \emph{IEEE Transactions on Information Theory} 18 (1972), 460--473.
\bibitem{blum2010} M. G. B. Blum. Approximate Bayesian computation: a nonparametric perspective. \emph{Journal of the American Statistical Association} 105 (2010), 1178--1187.
\bibitem{blum2013} M. G. B. Blum, M. A. Nunes, D. Prangle and S. A. Sisson. A comparative review of dimension reduction methods in approximate Bayesian computation. \emph{Statistical Science} 28 (2013), 189--208.
\bibitem{bobkovledoux2019} S. Bobkov and M. Ledoux. One dimensional empirical measures, order statistics, and Kantorovich transport distances. \emph{Memoirs of the American Mathematical Society} 261 (2019), 1--126.
\bibitem{bobylevcercignani1999} A. V. Bobylev and C. Cercignani. On the rate of entropy production for the Boltzmann equation. \emph{Journal of Statistical Physics} 94 (1999), 603--618.
\bibitem{bgs2016} T. Bodineau, I. Gallagher and L. Saint Raymond. The Brownian motion as the limit of a deterministic system of hard spheres. \emph{Inventiones Mathematicae} 203 (2016), 493--553.
\bibitem{bgss2023} T. Bodineau, I. Gallagher, L. Saint Raymond and S. Simonella. Statistical dynamics of a hard sphere gas: fluctuating Boltzmann equation and large deviations. \emph{Annals of Mathematics} 198 (2023), 1047--1201.
\bibitem{bgss2023cpam} T. Bodineau, I. Gallagher, L. Saint Raymond and S. Simonella. Long time correlations for a hard sphere gas at equilibrium. \emph{Communications on Pure and Applied Mathematics} 76 (2023), 3852--3911.
\bibitem{bgss2024} T. Bodineau, I. Gallagher, L. Saint Raymond and S. Simonella. Long time derivation at equilibrium of the fluctuating Boltzmann equation. \emph{The Annals of Probability} 52 (2024), 217--295.
\bibitem{bgss2026} T. Bodineau, I. Gallagher, L. Saint Raymond and S. Simonella. Derivation of the Boltzmann equation from hard sphere dynamics (after Y. Deng, Z. Hani and X. Ma). \emph{S\'eminaire Bourbaki}, expos\'e 1247 (January 2026). Preprint, arXiv:2602.04407, 2026.
\bibitem{bogachev1998} V. I. Bogachev. \emph{Gaussian measures}. Mathematical Surveys and Monographs 62, American Mathematical Society, 1998.
\bibitem{bohrnickl2024} J. Bohr and R. Nickl. On log concave approximations of high dimensional posterior measures and stability properties in nonlinear inverse problems. \emph{Annales de l'Institut Henri Poincar\'e, Probabilit\'es et Statistiques} 60 (2024), 2619--2667.
\bibitem{bbs1983} C. Boldrighini, L. A. Bunimovich and Ya. G. Sinai. On the Boltzmann equation for the Lorentz gas. \emph{Journal of Statistical Physics} 32 (1983), 477--501.
\bibitem{bollerslev1986} T. Bollerslev. Generalized autoregressive conditional heteroskedasticity. \emph{Journal of Econometrics} 31 (1986), 307--327.
\bibitem{boltzmann1868} L. Boltzmann. Studien \"uber das Gleichgewicht der lebendigen Kraft zwischen bewegten materiellen Punkten. \emph{Sitzungsberichte der Akademie der Wissenschaften Wien} 58 (1868), 517--560.
\bibitem{boltzmann1871} L. Boltzmann. Einige allgemeine S\"atze \"uber W\"armegleichgewicht. \emph{Sitzungsberichte der Akademie der Wissenschaften Wien} 63 (1871), 679--711.
\bibitem{boltzmann1872} L. Boltzmann. Weitere Studien \"uber das W\"armegleichgewicht unter Gasmolek\"ulen. \emph{Sitzungsberichte der Akademie der Wissenschaften Wien} 66 (1872), 275--370.
\bibitem{boltzmann1877} L. Boltzmann. \"Uber die Beziehung zwischen dem zweiten Hauptsatze der mechanischen W\"armetheorie und der Wahrscheinlichkeitsrechnung respektive den S\"atzen \"uber das W\"armegleichgewicht. \emph{Sitzungsberichte der Akademie der Wissenschaften Wien} 76 (1877), 373--435.
\bibitem{boltzmann1896} L. Boltzmann. \emph{Vorlesungen \"uber Gastheorie}. Two volumes, J. A. Barth, Leipzig, 1896 and 1898. English translation by S. G. Brush: \emph{Lectures on gas theory}, University of California Press, 1964.
\bibitem{bondesson1992} L. Bondesson. \emph{Generalized gamma convolutions and related classes of distributions and densities}. Lecture Notes in Statistics 76, Springer, 1992.
\bibitem{bosch2016} P. Bosch and T. Simon. A proof of Bondesson's conjecture on stable densities. \emph{Arkiv f\"or Matematik} 54 (2016), 31--38.
\bibitem{bougerol1983} P. Bougerol. Exemples de th\'eor\`emes locaux sur les groupes r\'esolubles. \emph{Annales de l'Institut Henri Poincar\'e, Probabilit\'es et Statistiques} 19 (1983), 369--391.
\bibitem{bramson1983} M. Bramson. Convergence of solutions of the Kolmogorov equation to travelling waves. \emph{Memoirs of the American Mathematical Society} 44, no. 285 (1983).
\bibitem{braunhepp1977} W. Braun and K. Hepp. The Vlasov dynamics and its fluctuations in the 1/N limit of interacting classical particles. \emph{Communications in Mathematical Physics} 56 (1977), 101--113.
\bibitem{braungriebel2009} J. Braun and M. Griebel. On a constructive proof of Kolmogorov's superposition theorem. \emph{Constructive Approximation} 30 (2009), 653--675.
\bibitem{breiman1957} L. Breiman. The individual ergodic theorem of information theory. \emph{The Annals of Mathematical Statistics} 28 (1957), 809--811; correction, 31 (1960), 809--810.
\bibitem{brenier1991} Y. Brenier. Polar factorization and monotone rearrangement of vector valued functions. \emph{Communications on Pure and Applied Mathematics} 44 (1991), 375--417.
\bibitem{brillinger2012} D. R. Brillinger. A generalized linear model with ``Gaussian'' regressor variables. In \emph{Selected Works of David Brillinger}, pages 589--606. Springer, 2012. Reprinted from \emph{A Festschrift for Erich L. Lehmann in honor of his sixty fifth birthday}, pages 97--114. Wadsworth, 1983.
\bibitem{brown1828} R. Brown. A brief account of microscopical observations made in the months of June, July and August 1827, on the particles contained in the pollen of plants; and on the general existence of active molecules in organic and inorganic bodies. \emph{Philosophical Magazine} 4 (1828), 161--173.
\bibitem{brown1959} R. G. Brown. \emph{Statistical forecasting for inventory control}. McGraw Hill, 1959.
\bibitem{brush1967} S. G. Brush. History of the Lenz Ising model. \emph{Reviews of Modern Physics} 39 (1967), 883--893.
\bibitem{bryc1995} W. Bryc. \emph{The normal distribution: characterizations with applications}. Lecture Notes in Statistics 100, Springer, 1995.
\bibitem{brynjarsdottirohagan2014} J. Brynjarsd\'ottir and A. O'Hagan. Learning about physical parameters: the importance of model discrepancy. \emph{Inverse Problems} 30 (2014), 114007.
\bibitem{bdsv2019} T. Buckmaster, C. De Lellis, L. Sz\'ekelyhidi Jr. and V. Vicol. Onsager's conjecture for admissible weak solutions. \emph{Communications on Pure and Applied Mathematics} 72 (2019), 229--274.
\bibitem{buckmastervicol2019} T. Buckmaster and V. Vicol. Nonuniqueness of weak solutions to the Navier Stokes equation. \emph{Annals of Mathematics} 189 (2019), 101--144.
\bibitem{buithanh2013} T. Bui Thanh, O. Ghattas, J. Martin and G. Stadler. A computational framework for infinite dimensional Bayesian inverse problems, Part I: the linearized case, with application to global seismic inversion. \emph{SIAM Journal on Scientific Computing} 35 (2013), A2494--A2523.
\bibitem{burgers1998} G. Burgers, P. J. van Leeuwen and G. Evensen. Analysis scheme in the ensemble Kalman filter. \emph{Monthly Weather Review} 126 (1998), 1719--1724.
\bibitem{ckn1982} L. Caffarelli, R. Kohn and L. Nirenberg. Partial regularity of suitable weak solutions of the Navier Stokes equations. \emph{Communications on Pure and Applied Mathematics} 35 (1982), 771--831.
\bibitem{caflisch1980} R. E. Caflisch. The fluid dynamic limit of the nonlinear Boltzmann equation. \emph{Communications on Pure and Applied Mathematics} 33 (1980), 651--666.
\bibitem{campbell1909} N. Campbell. The study of discontinuous phenomena. \emph{Proceedings of the Cambridge Philosophical Society} 15 (1909), 117--136.
\bibitem{caratheodory1919} C. Carath\'eodory. \"Uber den Wiederkehrsatz von Poincar\'e. \emph{Sitzungsberichte der Preussischen Akademie der Wissenschaften} (1919), 580--584.
\bibitem{carleman1933} T. Carleman. Sur la th\'eorie de l'\'equation int\'egrodiff\'erentielle de Boltzmann. \emph{Acta Mathematica} 60 (1933), 91--146.
\bibitem{carlen1984} E. A. Carlen. Conservative diffusions. \emph{Communications in Mathematical Physics} 94 (1984), 293--315.
\bibitem{carlen1991} E. A. Carlen. Superadditivity of Fisher's information and logarithmic Sobolev inequalities. \emph{Journal of Functional Analysis} 101 (1991), 194--211.
\bibitem{ccl2003} E. A. Carlen, M. C. Carvalho and M. Loss. Determination of the spectral gap for Kac's master equation and related stochastic evolution. \emph{Acta Mathematica} 191 (2003), 1--54.
\bibitem{cclv2010} E. A. Carlen, M. C. Carvalho, J. Le Roux, M. Loss and C. Villani. Entropy and chaos in the Kac model. \emph{Kinetic and Related Models} 3 (2010), 85--122.
\bibitem{carlencarvalhoeinav2018} E. A. Carlen, M. C. Carvalho and A. Einav. Entropy production inequalities for the Kac walk. \emph{Kinetic and Related Models} 11 (2018), 219--238.
\bibitem{carlier2016} G. Carlier, V. Chernozhukov and A. Galichon. Vector quantile regression: an optimal transport approach. \emph{The Annals of Statistics} 44 (2016), 1165--1192.
\bibitem{cps1992} B. P. Carlin, N. G. Polson and D. S. Stoffer. A Monte Carlo approach to nonnormal and nonlinear state space modeling. \emph{Journal of the American Statistical Association} 87 (1992), 493--500.
\bibitem{carvalho2010} C. M. Carvalho, N. G. Polson and J. G. Scott. The horseshoe estimator for sparse signals. \emph{Biometrika} 97 (2010), 465--480.
\bibitem{cjlp2010} C. M. Carvalho, M. S. Johannes, H. F. Lopes and N. G. Polson. Particle learning and smoothing. \emph{Statistical Science} 25 (2010), 88--106.
\bibitem{clpt2010} C. M. Carvalho, H. F. Lopes, N. G. Polson and M. A. Taddy. Particle learning for general mixtures. \emph{Bayesian Analysis} 5 (2010), 709--740.
\bibitem{castillonickl2013} I. Castillo and R. Nickl. Nonparametric Bernstein von Mises theorems in Gaussian white noise. \emph{The Annals of Statistics} 41 (2013), 1999--2028.
\bibitem{castrenickl2026} A. Castre and R. Nickl. On gradient stability in nonlinear PDE models and inference in interacting particle systems. Preprint, arXiv:2601.10326, 2026.
\bibitem{cip1994} C. Cercignani, R. Illner and M. Pulvirenti. \emph{The mathematical theory of dilute gases}. Springer, 1994.
\bibitem{cdg2011} F. C\'erou, P. Del Moral and A. Guyader. A nonasymptotic theorem for unnormalized Feynman Kac particle models. \emph{Annales de l'Institut Henri Poincar\'e, Probabilit\'es et Statistiques} 47 (2011), 629--649.
\bibitem{chapman1916} S. Chapman. On the law of distribution of molecular velocities, and on the theory of viscosity and thermal conduction, in a nonuniform simple monatomic gas. \emph{Philosophical Transactions of the Royal Society A} 216 (1916), 279--348.
\bibitem{chaudhari2018} P. Chaudhari and S. Soatto. Stochastic gradient descent performs variational inference, converges to limit cycles for deep networks. In \emph{2018 Information Theory and Applications Workshop}, pages 1--10. IEEE, 2018.
\bibitem{chenliu2000} R. Chen and J. S. Liu. Mixture Kalman filters. \emph{Journal of the Royal Statistical Society B} 62 (2000), 493--508.
\bibitem{chen2021} Y. Chen, B. Hosseini, H. Owhadi and A. M. Stuart. Solving and learning nonlinear PDEs with Gaussian processes. \emph{Journal of Computational Physics} 447 (2021), 110668.
\bibitem{chenhou2022} J. Chen and T. Y. Hou. Stable nearly self similar blowup of the 2D Boussinesq and 3D Euler equations with smooth data I: analysis. Preprint, arXiv:2210.07191, 2022.
\bibitem{chenhou2025} J. Chen and T. Y. Hou. Stable nearly self similar blowup of the 2D Boussinesq and 3D Euler equations with smooth data II: rigorous numerics. \emph{Multiscale Modeling \& Simulation} 23 (2025), 25--130.
\bibitem{chernoff1952} H. Chernoff. A measure of asymptotic efficiency for tests of a hypothesis based on the sum of observations. \emph{The Annals of Mathematical Statistics} 23 (1952), 493--507.
\bibitem{chernoff1956} H. Chernoff. Large sample theory: parametric case. \emph{The Annals of Mathematical Statistics} 27 (1956), 1--22.
\bibitem{chernoff1981} H. Chernoff. A note on an inequality involving the normal distribution. \emph{The Annals of Probability} 9 (1981), 533--535.
\bibitem{chernozhukov2018} V. Chernozhukov, D. Chetverikov, M. Demirer, E. Duflo, C. Hansen, W. Newey and J. Robins. Double/debiased machine learning for treatment and structural parameters. \emph{The Econometrics Journal} 21 (2018), C1--C68.
\bibitem{chewi2025} S. Chewi, M. A. Erdogdu, M. Li, R. Shen and M. S. Zhang. Analysis of Langevin Monte Carlo from Poincar\'e to log Sobolev. \emph{Foundations of Computational Mathematics} 25 (2025), 1345--1395.
\bibitem{chirikov1979} B. V. Chirikov. A universal instability of many dimensional oscillator systems. \emph{Physics Reports} 52 (1979), 263--379.
\bibitem{chopin2010} N. Chopin, A. Iacobucci, J. M. Marin, K. Mengersen, C. P. Robert, R. Ryder and C. Sch\"afer. On particle learning. Preprint, arXiv:1006.0554, 2010.
\bibitem{christenfox2005} J. A. Christen and C. Fox. Markov chain Monte Carlo using an approximation. \emph{Journal of Computational and Graphical Statistics} 14 (2005), 795--810.
\bibitem{ciesielski1962} Z. Ciesielski and S. J. Taylor. First passage times and sojourn times for Brownian motion in space and the exact Hausdorff measure of the sample path. \emph{Transactions of the American Mathematical Society} 103 (1962), 434--450.
\bibitem{clarkebarron1990} B. S. Clarke and A. R. Barron. Information theoretic asymptotics of Bayes methods. \emph{IEEE Transactions on Information Theory} 36 (1990), 453--471.
\bibitem{clarkebarron1994} B. S. Clarke and A. R. Barron. Jeffreys' prior is asymptotically least favorable under entropy risk. \emph{Journal of Statistical Planning and Inference} 41 (1994), 37--60.
\bibitem{clauser1969} J. F. Clauser, M. A. Horne, A. Shimony and R. A. Holt. Proposed experiment to test local hidden variable theories. \emph{Physical Review Letters} 23 (1969), 880--884.
\bibitem{clausius1865} R. Clausius. Ueber verschiedene f\"ur die Anwendung bequeme Formen der Hauptgleichungen der mechanischen W\"armetheorie. \emph{Annalen der Physik und Chemie} 125 (1865), 353--400.
\bibitem{clement2002} E. Cl\'ement, D. Lamberton and P. Protter. An analysis of a least squares regression method for American option pricing. \emph{Finance and Stochastics} 6 (2002), 449--471.
\bibitem{cole1951} J. D. Cole. On a quasilinear parabolic equation occurring in aerodynamics. \emph{Quarterly of Applied Mathematics} 9 (1951), 225--236.
\bibitem{cet1994} P. Constantin, W. E and E. S. Titi. Onsager's conjecture on the energy conservation for solutions of Euler's equation. \emph{Communications in Mathematical Physics} 165 (1994), 207--209.
\bibitem{corderoerausquin2001} D. Cordero Erausquin, R. J. McCann and M. Schmuckenschl\"ager. A Riemannian interpolation inequality \`a la Borell, Brascamp and Lieb. \emph{Inventiones Mathematicae} 146 (2001), 219--257.
\bibitem{cordoba2023} D. C\'ordoba and L. Mart\'inez Zoroa. Blow up for the incompressible 3D Euler equations with uniform $C^{1,1/2-\varepsilon}\cap L^2$ force. Preprint, arXiv:2309.08495, 2023.
\bibitem{cordoba2024} D. C\'ordoba and L. Mart\'inez Zoroa. Finite time singularities of smooth solutions for the 2D incompressible porous media (IPM) equation with a smooth source. Preprint, arXiv:2410.22920, 2024.
\bibitem{cordoba2026} D. C\'ordoba, L. Mart\'inez Zoroa and F. Zheng. Finite time blow up for the hypodissipative Navier Stokes equations with a force in $L^1_tC^{1,\varepsilon}_x\cap L^\infty_tL^2_x$. \emph{Archive for Rational Mechanics and Analysis} 250 (2026), 38.
\bibitem{corry2004} L. Corry. \emph{David Hilbert and the axiomatization of physics (1898 to 1918): from Grundlagen der Geometrie to Grundlagen der Physik}. Kluwer, 2004.
\bibitem{cotter2013} S. L. Cotter, G. O. Roberts, A. M. Stuart and D. White. MCMC methods for functions: modifying old algorithms to make them faster. \emph{Statistical Science} 28 (2013), 424--446.
\bibitem{courtault2000} J. M. Courtault, Y. Kabanov, B. Bru, P. Cr\'epel, I. Lebon and A. Le Marchand. Louis Bachelier on the centenary of Th\'eorie de la sp\'eculation. \emph{Mathematical Finance} 10 (2000), 339--353.
\bibitem{covertomas} T. M. Cover and J. A. Thomas. \emph{Elements of information theory}. Second edition, Wiley, 2006.
\bibitem{cox1946} R. T. Cox. Probability, frequency and reasonable expectation. \emph{American Journal of Physics} 14 (1946), 1--13.
\bibitem{cranmer2020} K. Cranmer, J. Brehmer and G. Louppe. The frontier of simulation based inference. \emph{Proceedings of the National Academy of Sciences} 117 (2020), 30055--30062.
\bibitem{crick1989} F. Crick. The recent excitement about neural networks. \emph{Nature} 337 (1989), 129--132.
\bibitem{crisandoucet2002} D. Crisan and A. Doucet. A survey of convergence results on particle filtering methods for practitioners. \emph{IEEE Transactions on Signal Processing} 50 (2002), 736--746.
\bibitem{crooks1999} G. E. Crooks. Entropy production fluctuation theorem and the nonequilibrium work relation for free energy differences. \emph{Physical Review E} 60 (1999), 2721--2726.
\bibitem{crooks2000} G. E. Crooks. Path ensemble averages in systems driven far from equilibrium. \emph{Physical Review E} 61 (2000), 2361--2366.
\bibitem{csiszar1975} I. Csisz\'ar. I divergence geometry of probability distributions and minimization problems. \emph{The Annals of Probability} 3 (1975), 146--158.
\bibitem{csiszar1984} I. Csisz\'ar. Sanov property, generalized I projection and a conditional limit theorem. \emph{The Annals of Probability} 12 (1984), 768--793.
\bibitem{csiszar1987} I. Csisz\'ar, T. M. Cover and B. S. Choi. Conditional limit theorems under Markov conditioning. \emph{IEEE Transactions on Information Theory} 33 (1987), 788--801.
\bibitem{csiszar1998} I. Csisz\'ar. The method of types. \emph{IEEE Transactions on Information Theory} 44 (1998), 2505--2523.
\bibitem{csiszarkorner2011} I. Csisz\'ar and J. K\"orner. \emph{Information theory: coding theorems for discrete memoryless systems}. Second edition, Cambridge University Press, 2011.
\bibitem{csv1994} G. Csordas, W. Smith and R. S. Varga. Lehmer pairs of zeros, the de Bruijn Newman constant $\Lambda$, and the Riemann hypothesis. \emph{Constructive Approximation} 10 (1994), 107--129.
\bibitem{cui2014} T. Cui, J. Martin, Y. M. Marzouk, A. Solonen and A. Spantini. Likelihood informed dimension reduction for nonlinear inverse problems. \emph{Inverse Problems} 30 (2014), 114015.
\bibitem{cuturi2013} M. Cuturi. Sinkhorn distances: lightspeed computation of optimal transport. \emph{Advances in Neural Information Processing Systems} 26 (2013), 2292--2300.
\bibitem{dalalyan2017} A. S. Dalalyan. Theoretical guarantees for approximate sampling from smooth and log concave densities. \emph{Journal of the Royal Statistical Society B} 79 (2017), 651--676.
\bibitem{damien1999} P. Damien, J. Wakefield and S. Walker. Gibbs sampling for Bayesian non conjugate and hierarchical models by using auxiliary variables. \emph{Journal of the Royal Statistical Society B} 61 (1999), 331--344.
\bibitem{darboux1904} G. Darboux, P. Appell and H. Poincar\'e. Examen critique des divers syst\`emes ou \'etudes graphologiques auxquels a donn\'e lieu le bordereau. Report to the Cour de cassation, 2 August 1904, reproduced in \emph{Journal \'Electronique d'Histoire des Probabilit\'es et de la Statistique} 1 (2005).
\bibitem{darmois1935} G. Darmois. Sur les lois de probabilit\'e \`a estimation exhaustive. \emph{Comptes Rendus de l'Acad\'emie des Sciences, Paris} 200 (1935), 1265--1266.
\bibitem{darmois1953} G. Darmois. Analyse g\'en\'erale des liaisons stochastiques: \'etude particuli\`ere de l'analyse factorielle lin\'eaire. \emph{Revue de l'Institut International de Statistique} 21 (1953), 2--8.
\bibitem{dashti2013} M. Dashti, K. J. H. Law, A. M. Stuart and J. Voss. MAP estimators and their consistency in Bayesian nonparametric inverse problems. \emph{Inverse Problems} 29 (2013), 095017.
\bibitem{dashti2017} M. Dashti and A. M. Stuart. The Bayesian approach to inverse problems. In \emph{Handbook of Uncertainty Quantification}, eds. R. Ghanem, D. Higdon and H. Owhadi, pages 311--428. Springer, 2017.
\bibitem{dattapolson2025} J. Datta and N. G. Polson. Polynomial log marginals and Tweedie's formula: when is Bayes possible? Preprint, arXiv:2509.05823, 2025.
\bibitem{davisson1981} L. Davisson, G. Longo and A. Sgarro. The error exponent for the noiseless encoding of finite ergodic Markov sources. \emph{IEEE Transactions on Information Theory} 27 (1981), 431--438.
\bibitem{debortoli2021} V. De Bortoli, J. Thornton, J. Heng and A. Doucet. Diffusion Schr\"odinger bridge with applications to score based generative modeling. \emph{Advances in Neural Information Processing Systems} 34 (2021), 17695--17709.
\bibitem{debruijn1950} N. G. de Bruijn. The roots of trigonometric integrals. \emph{Duke Mathematical Journal} 17 (1950), 197--226.
\bibitem{definetti1937} B. de Finetti. La pr\'evision: ses lois logiques, ses sources subjectives. \emph{Annales de l'Institut Henri Poincar\'e} 7 (1937), 1--68.
\bibitem{degiorgi1957} E. De Giorgi. Sulla differenziabilit\`a e l'analiticit\`a delle estremali degli integrali multipli regolari. \emph{Memorie della Accademia delle Scienze di Torino, Classe di Scienze Fisiche, Matematiche e Naturali} (3) 3 (1957), 25--43.
\bibitem{delellisszekelyhidi2009} C. De Lellis and L. Sz\'ekelyhidi. The Euler equations as a differential inclusion. \emph{Annals of Mathematics} 170 (2009), 1417--1436.
\bibitem{delellisszekelyhidi2017} C. De Lellis and L. Sz\'ekelyhidi. High dimensionality and $h$ principle in PDE. \emph{Bulletin of the American Mathematical Society} 54 (2017), 247--282.
\bibitem{delmoral1998} P. Del Moral. Measure valued processes and interacting particle systems. Application to nonlinear filtering problems. \emph{The Annals of Applied Probability} 8 (1998), 438--495.
\bibitem{delmoral2001} P. Del Moral, J. Jacod and P. Protter. The Monte Carlo method for filtering with discrete time observations. \emph{Probability Theory and Related Fields} 120 (2001), 346--368.
\bibitem{delmoralguionnet2001} P. Del Moral and A. Guionnet. On the stability of interacting processes with applications to filtering and genetic algorithms. \emph{Annales de l'Institut Henri Poincar\'e, Probabilit\'es et Statistiques} 37 (2001), 155--194.
\bibitem{delmoral2004} P. Del Moral. \emph{Feynman Kac formulae: genealogical and interacting particle systems with applications}. Springer, 2004.
\bibitem{delmoral2006} P. Del Moral, A. Doucet and A. Jasra. Sequential Monte Carlo samplers. \emph{Journal of the Royal Statistical Society B} 68 (2006), 411--436.
\bibitem{dembozeitouni} A. Dembo and O. Zeitouni. \emph{Large deviations techniques and applications}. Second edition, Springer, 1998.
\bibitem{demingstephan1940} W. E. Deming and F. F. Stephan. On a least squares adjustment of a sampled frequency table when the expected marginal totals are known. \emph{The Annals of Mathematical Statistics} 11 (1940), 427--444.
\bibitem{denghani2023} Y. Deng and Z. Hani. Full derivation of the wave kinetic equation. \emph{Inventiones Mathematicae} 233 (2023), 543--724.
\bibitem{denghanima2024} Y. Deng, Z. Hani and X. Ma. Long time derivation of the Boltzmann equation from hard sphere dynamics. Preprint, arXiv:2408.07818, 2024.
\bibitem{denghanima2025} Y. Deng, Z. Hani and X. Ma. Hilbert's sixth problem: derivation of fluid equations via Boltzmann's kinetic theory. Preprint, arXiv:2503.01800, 2025.
\bibitem{deshpande2015} Y. Deshpande and A. Montanari. Finding hidden cliques of size $\sqrt{N/e}$ in nearly linear time. \emph{Foundations of Computational Mathematics} 15 (2015), 1069--1128.
\bibitem{desvillettesvillani2005} L. Desvillettes and C. Villani. On the trend to global equilibrium for spatially inhomogeneous kinetic systems: the Boltzmann equation. \emph{Inventiones Mathematicae} 159 (2005), 245--316.
\bibitem{diaconisfreedman1980} P. Diaconis and D. Freedman. Finite exchangeable sequences. \emph{The Annals of Probability} 8 (1980), 745--764.
\bibitem{diaconisshahshahani1981} P. Diaconis and M. Shahshahani. Generating a random permutation with random transpositions. \emph{Zeitschrift f\"ur Wahrscheinlichkeitstheorie und verwandte Gebiete} 57 (1981), 159--179.
\bibitem{diaconisfreedman1986} P. Diaconis and D. Freedman. On the consistency of Bayes estimates. \emph{The Annals of Statistics} 14 (1986), 1--26.
\bibitem{diaconisfreedman1987} P. Diaconis and D. Freedman. A dozen de Finetti style results in search of a theory. \emph{Annales de l'Institut Henri Poincar\'e, Probabilit\'es et Statistiques} 23 (1987), 397--423.
\bibitem{diaconis1988} P. Diaconis. \emph{Group representations in probability and statistics}. Institute of Mathematical Statistics Lecture Notes Monograph Series 11, Institute of Mathematical Statistics, 1988.
\bibitem{diaconisfreedman1988} P. Diaconis and D. A. Freedman. Conditional limit theorems for exponential families and finite versions of de Finetti's theorem. \emph{Journal of Theoretical Probability} 1 (1988), 381--410.
\bibitem{diaconis2007coin} P. Diaconis, S. Holmes and R. Montgomery. Dynamical bias in the coin toss. \emph{SIAM Review} 49 (2007), 211--235.
\bibitem{diperna1989} R. J. DiPerna and P. L. Lions. On the Cauchy problem for Boltzmann equations: global existence and weak stability. \emph{Annals of Mathematics} 130 (1989), 321--366.
\bibitem{dobrushin1979} R. L. Dobrushin. Vlasov equations. \emph{Functional Analysis and Its Applications} 13 (1979), 115--123.
\bibitem{dodwell2015} T. J. Dodwell, C. Ketelsen, R. Scheichl and A. L. Teckentrup. A hierarchical multilevel Markov chain Monte Carlo algorithm with applications to uncertainty quantification in subsurface flow. \emph{SIAM/ASA Journal on Uncertainty Quantification} 3 (2015), 1075--1108.
\bibitem{donoho2009} D. L. Donoho, A. Maleki and A. Montanari. Message passing algorithms for compressed sensing. \emph{Proceedings of the National Academy of Sciences} 106 (2009), 18914--18919.
\bibitem{donsker1975} M. D. Donsker and S. R. S. Varadhan. Asymptotic evaluation of certain Markov process expectations for large time, I. \emph{Communications on Pure and Applied Mathematics} 28 (1975), 1--47.
\bibitem{doob1949} J. L. Doob. Application of the theory of martingales. In \emph{Actes du Colloque International Le Calcul des Probabilit\'es et ses Applications} (Lyon, 1948), pages 23--27. CNRS, Paris, 1949.
\bibitem{dga2000} A. Doucet, S. Godsill and C. Andrieu. On sequential Monte Carlo sampling methods for Bayesian filtering. \emph{Statistics and Computing} 10 (2000), 197--208.
\bibitem{duane1987} S. Duane, A. D. Kennedy, B. J. Pendleton and D. Roweth. Hybrid Monte Carlo. \emph{Physics Letters B} 195 (1987), 216--222.
\bibitem{duchonrobert2000} J. Duchon and R. Robert. Inertial energy dissipation for weak solutions of incompressible Euler and Navier Stokes equations. \emph{Nonlinearity} 13 (2000), 249--255.
\bibitem{dudley2002} R. M. Dudley. \emph{Real analysis and probability}. Second edition, Cambridge Studies in Advanced Mathematics 74, Cambridge University Press, 2002.
\bibitem{dufresne1990} D. Dufresne. The distribution of a perpetuity, with applications to risk theory and pension funding. \emph{Scandinavian Actuarial Journal} (1990), 39--79.
\bibitem{durkalski2003} V. L. Durkalski, Y. Y. Palesch, S. R. Lipsitz and P. F. Rust. Analysis of clustered matched pair data. \emph{Statistics in Medicine} 22 (2003), 2417--2428.
\bibitem{durmusmoulines2017} A. Durmus and \'E. Moulines. Nonasymptotic convergence analysis for the unadjusted Langevin algorithm. \emph{The Annals of Applied Probability} 27 (2017), 1551--1587.
\bibitem{durmusmoulines2019} A. Durmus and \'E. Moulines. High dimensional Bayesian inference via the unadjusted Langevin algorithm. \emph{Bernoulli} 25 (2019), 2854--2882.
\bibitem{dyer1991} M. Dyer, A. Frieze and R. Kannan. A random polynomial time algorithm for approximating the volume of convex bodies. \emph{Journal of the ACM} 38 (1991), 1--17.
\bibitem{dyson1962} F. J. Dyson. A Brownian motion model for the eigenvalues of a random matrix. \emph{Journal of Mathematical Physics} 3 (1962), 1191--1198.
\bibitem{dysonfalk1992} F. J. Dyson and H. Falk. Period of a discrete cat mapping. \emph{American Mathematical Monthly} 99 (1992), 603--614.
\bibitem{borel1906} \'E. Borel. Sur les principes de la th\'eorie cin\'etique des gaz. \emph{Annales scientifiques de l'\'Ecole Normale Sup\'erieure} 23 (1906), 9--32.
\bibitem{borel1909} \'E. Borel. Les probabilit\'es d\'enombrables et leurs applications arithm\'etiques. \emph{Rendiconti del Circolo Matematico di Palermo} 27 (1909), 247--271.
\bibitem{borel1909elements} \'E. Borel. \emph{\'El\'ements de la th\'eorie des probabilit\'es}. A. Hermann, Paris, 1909.
\bibitem{borel1914} \'E. Borel. \emph{Introduction g\'eom\'etrique \`a quelques th\'eories physiques}. Gauthier Villars, Paris, 1914.
\bibitem{ekms2000} W. E, K. Khanin, A. Mazel and Ya. Sinai. Invariant measures for Burgers equation with stochastic forcing. \emph{Annals of Mathematics} 151 (2000), 877--960.
\bibitem{evandeneijnden2000} W. E and E. Vanden Eijnden. Statistical theory for the stochastic Burgers equation in the inviscid limit. \emph{Communications on Pure and Applied Mathematics} 53 (2000), 852--901.
\bibitem{ebert2000} U. Ebert and W. van Saarloos. Front propagation into unstable states: universal algebraic convergence towards uniformly translating pulled fronts. \emph{Physica D} 146 (2000), 1--99.
\bibitem{edwards1988} R. G. Edwards and A. D. Sokal. Generalization of the Fortuin Kasteleyn Swendsen Wang representation and Monte Carlo algorithm. \emph{Physical Review D} 38 (1988), 2009--2012.
\bibitem{efron2011} B. Efron. Tweedie's formula and selection bias. \emph{Journal of the American Statistical Association} 106 (2011), 1602--1614.
\bibitem{ehrenfest1911quanta} P. Ehrenfest. Welche Z\"uge der Lichtquantenhypothese spielen in der Theorie der W\"armestrahlung eine wesentliche Rolle? \emph{Annalen der Physik} 36 (1911), 91--118.
\bibitem{ehrenfest1911} P. Ehrenfest and T. Ehrenfest. Begriffliche Grundlagen der statistischen Auffassung in der Mechanik. In \emph{Encyklop\"adie der mathematischen Wissenschaften}, volume IV.2.II, number 6, pages 773--860. Teubner, 1911. English translation: \emph{The conceptual foundations of the statistical approach in mechanics}, Cornell University Press, 1959.
\bibitem{einav2011} A. Einav. On Villani's conjecture concerning entropy production for the Kac master equation. \emph{Kinetic and Related Models} 4 (2011), 479--497.
\bibitem{einstein1905} A. Einstein. \"Uber die von der molekularkinetischen Theorie der W\"arme geforderte Bewegung von in ruhenden Fl\"ussigkeiten suspendierten Teilchen. \emph{Annalen der Physik} 17 (1905), 549--560.
\bibitem{elerian2001} O. Elerian, S. Chib and N. Shephard. Likelihood inference for discretely observed nonlinear diffusions. \emph{Econometrica} 69 (2001), 959--993.
\bibitem{elgindi2021} T. M. Elgindi. Finite time singularity formation for $C^{1,\alpha}$ solutions to the incompressible Euler equations on $\mathbb{R}^3$. \emph{Annals of Mathematics} 194 (2021), 647--727.
\bibitem{elmoselhymarzouk2012} T. A. El Moselhy and Y. M. Marzouk. Bayesian inference with optimal maps. \emph{Journal of Computational Physics} 231 (2012), 7815--7850.
\bibitem{engel1992} E. M. R. A. Engel. \emph{A road to randomness in physical systems}. Lecture Notes in Statistics 71, Springer, 1992.
\bibitem{engle1982} R. F. Engle. Autoregressive conditional heteroscedasticity with estimates of the variance of United Kingdom inflation. \emph{Econometrica} 50 (1982), 987--1007.
\bibitem{enskog1917} D. Enskog. Kinetische Theorie der Vorg\"ange in m\"assig verd\"unnten Gasen. I. Allgemeiner Teil. Dissertation, Almqvist \& Wiksell, Uppsala, 1917.
\bibitem{eraker2001} B. Eraker. MCMC analysis of diffusion models with application to finance. \emph{Journal of Business and Economic Statistics} 19 (2001), 177--191.
\bibitem{ernst2015} O. G. Ernst, B. Sprungk and H. J. Starkloff. Analysis of the ensemble and polynomial chaos Kalman filters in Bayesian inverse problems. \emph{SIAM/ASA Journal on Uncertainty Quantification} 3 (2015), 823--851.
\bibitem{ess2003} L. Escauriaza, G. A. Seregin and V. \v{S}ver\'ak. $L_{3,\infty}$ solutions of the Navier Stokes equations and backward uniqueness. \emph{Russian Mathematical Surveys} 58 (2003), 211--250.
\bibitem{ethier1986} S. N. Ethier and T. G. Kurtz. \emph{Markov processes: characterization and convergence}. Wiley, 1986.
\bibitem{evensen1994} G. Evensen. Sequential data assimilation with a nonlinear quasigeostrophic model using Monte Carlo methods to forecast error statistics. \emph{Journal of Geophysical Research} 99 (1994), 10143--10162.
\bibitem{eyink2003} G. L. Eyink. Local 4/5 law and energy dissipation anomaly in turbulence. \emph{Nonlinearity} 16 (2003), 137--145.
\bibitem{eyinksreenivasan2006} G. L. Eyink and K. R. Sreenivasan. Onsager and the theory of hydrodynamic turbulence. \emph{Reviews of Modern Physics} 78 (2006), 87--135.
\bibitem{fearnhead2002} P. Fearnhead. Markov chain Monte Carlo, sufficient statistics, and particle filters. \emph{Journal of Computational and Graphical Statistics} 11 (2002), 848--862.
\bibitem{fearnheadprangle2012} P. Fearnhead and D. Prangle. Constructing summary statistics for approximate Bayesian computation: semiautomatic approximate Bayesian computation (with discussion). \emph{Journal of the Royal Statistical Society B} 74 (2012), 419--474.
\bibitem{ferguson1973} T. S. Ferguson. A Bayesian analysis of some nonparametric problems. \emph{The Annals of Statistics} 1 (1973), 209--230.
\bibitem{fieller1954} E. C. Fieller. Some problems in interval estimation. \emph{Journal of the Royal Statistical Society B} 16 (1954), 175--185.
\bibitem{fine1982} A. Fine. Hidden variables, joint probability, and the Bell inequalities. \emph{Physical Review Letters} 48 (1982), 291--295.
\bibitem{finocchioschmidthieber2023} G. Finocchio and J. Schmidt Hieber. Posterior contraction for deep Gaussian process priors. \emph{Journal of Machine Learning Research} 24 (2023), 1--49.
\bibitem{fisher1925} R. A. Fisher. Theory of statistical estimation. \emph{Proceedings of the Cambridge Philosophical Society} 22 (1925), 700--725.
\bibitem{fisher1937} R. A. Fisher. The wave of advance of advantageous genes. \emph{Annals of Eugenics} 7 (1937), 355--369.
\bibitem{follmer1985} H. F\"ollmer. An entropy approach to the time reversal of diffusion processes. In \emph{Stochastic differential systems filtering and control}, Lecture Notes in Control and Information Sciences 69, pages 156--163. Springer, 1985.
\bibitem{follmer1988} H. F\"ollmer. Random fields and diffusion processes. In \emph{\'Ecole d'\'Et\'e de Probabilit\'es de Saint Flour XV to XVII, 1985 to 1987}, pages 101--203. Lecture Notes in Mathematics 1362, Springer, 1988.
\bibitem{fongholmeswalker2023} E. Fong, C. Holmes and S. G. Walker. Martingale posterior distributions. \emph{Journal of the Royal Statistical Society B} 85 (2023), 1357--1391.
\bibitem{fortet1940} R. Fortet. R\'esolution d'un syst\`eme d'\'equations de M. Schr\"odinger. \emph{Journal de Math\'ematiques Pures et Appliqu\'ees} 19 (1940), 83--105.
\bibitem{franklinlorenz1989} J. Franklin and J. Lorenz. On the scaling of multidimensional matrices. \emph{Linear Algebra and its Applications} 114/115 (1989), 717--735.
\bibitem{frechet1921} M. Fr\'echet. Remarque sur les probabilit\'es continues. \emph{Bulletin des sciences math\'ematiques} 45 (1921), 87--88.
\bibitem{freedman1963} D. A. Freedman. On the asymptotic behavior of Bayes' estimates in the discrete case. \emph{The Annals of Mathematical Statistics} 34 (1963), 1386--1403.
\bibitem{friedli2017} S. Friedli and Y. Velenik. \emph{Statistical mechanics of lattice systems: a concrete mathematical introduction}. Cambridge University Press, 2017.
\bibitem{friedman1975} A. Friedman. \emph{Stochastic differential equations and applications}, Volume 1. Academic Press, 1975.
\bibitem{frieze1994} A. Frieze, R. Kannan and N. Polson. Sampling from log concave distributions. \emph{The Annals of Applied Probability} 4 (1994), 812--837.
\bibitem{frisch1995} U. Frisch. \emph{Turbulence: the legacy of A. N. Kolmogorov}. Cambridge University Press, 1995.
\bibitem{gallager1963} R. G. Gallager. \emph{Low density parity check codes}. MIT Press, 1963.
\bibitem{gst2013} I. Gallagher, L. Saint Raymond and B. Texier. \emph{From Newton to Boltzmann: hard spheres and short range potentials}. European Mathematical Society, 2014.
\bibitem{gallagher2019} I. Gallagher. From Newton to Navier Stokes, or how to connect fluid mechanics equations from microscopic to macroscopic scales. \emph{Bulletin of the American Mathematical Society} 56 (2019), 65--85.
\bibitem{gallaghertristani2020} I. Gallagher and I. Tristani. On the convergence of smooth solutions from Boltzmann to Navier Stokes. \emph{Annales Henri Lebesgue} 3 (2020), 561--614.
\bibitem{gallavotti1969} G. Gallavotti. Divergences and the approach to equilibrium in the Lorentz and the wind tree models. \emph{Physical Review} 185 (1969), 308--322.
\bibitem{gallavottiornstein1974} G. Gallavotti and D. S. Ornstein. Billiards and Bernoulli schemes. \emph{Communications in Mathematical Physics} 38 (1974), 83--101.
\bibitem{gamarnik2021} D. Gamarnik. The overlap gap property: a topological barrier to optimizing over random structures. \emph{Proceedings of the National Academy of Sciences} 118 (2021), e2108492118.
\bibitem{gaowangxia2026} F. Gao, S. Wang and X. Xia. Villani's conjecture for Kac's walk. Preprint, arXiv:2609.22675, 2026.
\bibitem{gardner1985} E. S. Gardner Jr. Exponential smoothing: the state of the art. \emph{Journal of Forecasting} 4 (1985), 1--28.
\bibitem{garsia1965} A. M. Garsia. A simple proof of E. Hopf's maximal ergodic theorem. \emph{Journal of Mathematics and Mechanics} 14 (1965), 381--382.
\bibitem{gauss1809} C. F. Gauss. \emph{Theoria motus corporum coelestium in sectionibus conicis solem ambientium}. Perthes and Besser, 1809.
\bibitem{genevay2019} A. Genevay, L. Chizat, F. Bach, M. Cuturi and G. Peyr\'e. Sample complexity of Sinkhorn divergences. In \emph{Proceedings of the 22nd International Conference on Artificial Intelligence and Statistics}, Proceedings of Machine Learning Research 89, pages 1574--1583, 2019.
\bibitem{georgii1995} H. O. Georgii. The equivalence of ensembles for classical systems of particles. \emph{Journal of Statistical Physics} 80 (1995), 1341--1378.
\bibitem{gerovitch2002} S. Gerovitch. \emph{From Newspeak to Cyberspeak: a history of Soviet cybernetics}. MIT Press, 2002.
\bibitem{gers2000} F. A. Gers, J. Schmidhuber and F. Cummins. Learning to forget: continual prediction with LSTM. \emph{Neural Computation} 12 (2000), 2451--2471.
\bibitem{gerstner2002} W. Gerstner and W. M. Kistler. \emph{Spiking neuron models: single neurons, populations, plasticity}. Cambridge University Press, 2002.
\bibitem{ghattaswillcox2021} O. Ghattas and K. Willcox. Learning physics based models from data: perspectives from inverse problems and model reduction. \emph{Acta Numerica} 30 (2021), 445--554.
\bibitem{gibbs1902} J. W. Gibbs. \emph{Elementary principles in statistical mechanics, developed with especial reference to the rational foundation of thermodynamics}. Charles Scribner's Sons, 1902.
\bibitem{ginenickl2016} E. Gin\'e and R. Nickl. \emph{Mathematical foundations of infinite dimensional statistical models}. Cambridge University Press, 2016.
\bibitem{giordanokekkonen2020} M. Giordano and H. Kekkonen. Bernstein von Mises theorems and uncertainty quantification for linear inverse problems. \emph{SIAM/ASA Journal on Uncertainty Quantification} 8 (2020), 342--373.
\bibitem{giordanonickl2020} M. Giordano and R. Nickl. Consistency of Bayesian inference with Gaussian process priors in an elliptic inverse problem. \emph{Inverse Problems} 36 (2020), 085001.
\bibitem{girosipoggio1989} F. Girosi and T. Poggio. Representation properties of networks: Kolmogorov's theorem is irrelevant. \emph{Neural Computation} 1 (1989), 465--469.
\bibitem{girsanov1960} I. V. Girsanov. On transforming a certain class of stochastic processes by absolutely continuous substitution of measures. \emph{Theory of Probability and Its Applications} 5 (1960), 285--301.
\bibitem{giustina2015} M. Giustina et al. Significant loophole free test of Bell's theorem with entangled photons. \emph{Physical Review Letters} 115 (2015), 250401.
\bibitem{glaeser1963} G. Glaeser. Racine carr\'ee d'une fonction diff\'erentiable. \emph{Annales de l'Institut Fourier} 13 (1963), 203--210.
\bibitem{gleason1957} A. M. Gleason. Measures on the closed subspaces of a Hilbert space. \emph{Journal of Mathematics and Mechanics} 6 (1957), 885--893.
\bibitem{glimmjaffe} J. Glimm and A. Jaffe. \emph{Quantum physics: a functional integral point of view}. Second edition, Springer, 1987.
\bibitem{goedel1931} K. G\"odel. \"Uber formal unentscheidbare S\"atze der Principia Mathematica und verwandter Systeme I. \emph{Monatshefte f\"ur Mathematik und Physik} 38 (1931), 173--198.
\bibitem{goedel1931b} K. G\"odel. Diskussion zur Grundlegung der Mathematik. \emph{Erkenntnis} 2 (1931), 147--151.
\bibitem{gdw2004} S. J. Godsill, A. Doucet and M. West. Monte Carlo smoothing for nonlinear time series. \emph{Journal of the American Statistical Association} 99 (2004), 156--168.
\bibitem{golselevermore2002} F. Golse and C. D. Levermore. Stokes Fourier and acoustic limits for the Boltzmann equation: convergence proofs. \emph{Communications on Pure and Applied Mathematics} 55 (2002), 336--393.
\bibitem{golsesaintraymond2004} F. Golse and L. Saint Raymond. The Navier Stokes limit of the Boltzmann equation for bounded collision kernels. \emph{Inventiones Mathematicae} 155 (2004), 81--161.
\bibitem{golse2008} F. Golse. On the periodic Lorentz gas and the Lorentz kinetic equation. \emph{Annales de la Facult\'e des Sciences de Toulouse, Math\'ematiques} 17 (2008), 735--749.
\bibitem{golsesaintraymond2009} F. Golse and L. Saint Raymond. The incompressible Navier Stokes limit of the Boltzmann equation for hard cutoff potentials. \emph{Journal de Math\'ematiques Pures et Appliqu\'ees} 91 (2009), 508--552.
\bibitem{golse2011} F. Golse. From the kinetic theory of gases to continuum mechanics. In \emph{Rarefied Gas Dynamics: 27th International Symposium}, AIP Conference Proceedings 1333, pages 15--27, 2011.
\bibitem{gorban2018} A. N. Gorban. Hilbert's sixth problem: the endless road to rigour. \emph{Philosophical Transactions of the Royal Society A} 376 (2018), 20170238.
\bibitem{gordin1969} M. I. Gordin. The central limit theorem for stationary processes. \emph{Doklady Akademii Nauk SSSR} 188 (1969), 739--741; English translation in \emph{Soviet Mathematics Doklady} 10 (1969), 1174--1176.
\bibitem{gss1993} N. J. Gordon, D. J. Salmond and A. F. M. Smith. Novel approach to nonlinear/non Gaussian Bayesian state estimation. \emph{IEE Proceedings F} 140 (1993), 107--113.
\bibitem{grad1949} H. Grad. On the kinetic theory of rarefied gases. \emph{Communications on Pure and Applied Mathematics} 2 (1949), 331--407.
\bibitem{grad1958} H. Grad. Principles of the kinetic theory of gases. In \emph{Handbuch der Physik}, vol. 12, pages 205--294. Springer, 1958.
\bibitem{grad1963} H. Grad. Asymptotic theory of the Boltzmann equation II. In \emph{Rarefied Gas Dynamics}, vol. 1, ed. J. A. Laurmann, pages 26--59. Academic Press, 1963.
\bibitem{gray2013} J. Gray. \emph{Henri Poincar\'e: a scientific biography}. Princeton University Press, 2013.
\bibitem{greenberg2019} D. S. Greenberg, M. Nonnenmacher and J. H. Macke. Automatic posterior transformation for likelihood free inference. In \emph{Proceedings of the 36th International Conference on Machine Learning}, PMLR 97, 2404--2414, 2019.
\bibitem{greene1979} J. M. Greene. A method for determining a stochastic transition. \emph{Journal of Mathematical Physics} 20 (1979), 1183--1201.
\bibitem{griffiths1970} R. B. Griffiths, C. A. Hurst and S. Sherman. Concavity of magnetization of an Ising ferromagnet in a positive external field. \emph{Journal of Mathematical Physics} 11 (1970), 790--795.
\bibitem{gross1975} L. Gross. Logarithmic Sobolev inequalities. \emph{American Journal of Mathematics} 97 (1975), 1061--1083.
\bibitem{gpv1988} M. Z. Guo, G. C. Papanicolaou and S. R. S. Varadhan. Nonlinear diffusion limit for a system with nearest neighbor interactions. \emph{Communications in Mathematical Physics} 118 (1988), 31--59.
\bibitem{guo2005} D. Guo, S. Shamai and S. Verd\'u. Mutual information and minimum mean square error in Gaussian channels. \emph{IEEE Transactions on Information Theory} 51 (2005), 1261--1282.
\bibitem{guojangjiang2009} Y. Guo, J. Jang and N. Jiang. Local Hilbert expansion for the Boltzmann equation. \emph{Kinetic and Related Models} 2 (2009), 205--214.
\bibitem{gyorfi2002} L. Gy\"orfi, M. Kohler, A. Krzy\.zak and H. Walk. \emph{A distribution free theory of nonparametric regression}. Springer, 2002.
\bibitem{hairer2014} M. Hairer, A. M. Stuart and S. J. Vollmer. Spectral gaps for a Metropolis Hastings algorithm in infinite dimensions. \emph{The Annals of Applied Probability} 24 (2014), 2455--2490.
\bibitem{halpern1999} J. Y. Halpern. A counter example to theorems of Cox and Fine. \emph{Journal of Artificial Intelligence Research} 10 (1999), 67--85.
\bibitem{hamilton1982} R. S. Hamilton. Three manifolds with positive Ricci curvature. \emph{Journal of Differential Geometry} 17 (1982), 255--306.
\bibitem{hamilton1989} J. D. Hamilton. A new approach to the economic analysis of nonstationary time series and the business cycle. \emph{Econometrica} 57 (1989), 357--384.
\bibitem{hmrt2022} T. Hastie, A. Montanari, S. Rosset and R. J. Tibshirani. Surprises in high dimensional ridgeless least squares interpolation. \emph{The Annals of Statistics} 50 (2022), 949--986.
\bibitem{hauraymischler2014} M. Hauray and S. Mischler. On Kac's chaos and related problems. \emph{Journal of Functional Analysis} 266 (2014), 6055--6157.
\bibitem{hebb1949} D. O. Hebb. \emph{The organization of behavior}. Wiley, 1949.
\bibitem{henonheiles1964} M. H\'enon and C. Heiles. The applicability of the third integral of motion: some numerical experiments. \emph{The Astronomical Journal} 69 (1964), 73--79.
\bibitem{hensen2015} B. Hensen et al. Loophole free Bell inequality violation using electron spins separated by 1.3 kilometres. \emph{Nature} 526 (2015), 682--686.
\bibitem{herschel1850} J. F. W. Herschel. Quetelet on probabilities. \emph{Edinburgh Review} 92 (1850), 1--57.
\bibitem{hewittsavage1955} E. Hewitt and L. J. Savage. Symmetric measures on Cartesian products. \emph{Transactions of the American Mathematical Society} 80 (1955), 470--501.
\bibitem{higdon1998} D. Higdon. A process convolution approach to modelling temperatures in the North Atlantic Ocean. \emph{Environmental and Ecological Statistics} 5 (1998), 173--190.
\bibitem{higdon1998b} D. Higdon. Auxiliary variable methods for Markov chain Monte Carlo with applications. \emph{Journal of the American Statistical Association} 93 (1998), 585--595.
\bibitem{higdon1999} D. Higdon, J. Swall and J. Kern. Nonstationary spatial modeling. In \emph{Bayesian Statistics 6}, pages 761--768. Oxford University Press, 1999.
\bibitem{higdon2002} D. Higdon. Space and space time modeling using process convolutions. In \emph{Quantitative Methods for Current Environmental Issues}, pages 37--56. Springer, 2002.
\bibitem{higdon2003} D. Higdon, H. Lee and C. Holloman. Markov chain Monte Carlo based approaches for inference in computationally intensive inverse problems. In \emph{Bayesian Statistics 7}, pages 181--197. Oxford University Press, 2003.
\bibitem{higdon2004} D. Higdon, M. Kennedy, J. C. Cavendish, J. A. Cafeo and R. D. Ryne. Combining field data and computer simulations for calibration and prediction. \emph{SIAM Journal on Scientific Computing} 26 (2004), 448--466.
\bibitem{higdon2008} D. Higdon, J. Gattiker, B. Williams and M. Rightley. Computer model calibration using high dimensional output. \emph{Journal of the American Statistical Association} 103 (2008), 570--583.
\bibitem{hilbert1895} D. Hilbert. Ueber die gerade Linie als k\"urzeste Verbindung zweier Punkte. \emph{Mathematische Annalen} 46 (1895), 91--96.
\bibitem{hilbert1900} D. Hilbert. Mathematische Probleme. \emph{Nachrichten von der K\"oniglichen Gesellschaft der Wissenschaften zu G\"ottingen} (1900), 253--297. English translation by M. W. Newson, \emph{Bulletin of the American Mathematical Society} 8 (1902), 437--479.
\bibitem{hilbert1912} D. Hilbert. Begr\"undung der kinetischen Gastheorie. \emph{Mathematische Annalen} 72 (1912), 562--577.
\bibitem{hilbert1915} D. Hilbert. Die Grundlagen der Physik. \emph{Nachrichten von der K\"oniglichen Gesellschaft der Wissenschaften zu G\"ottingen} (1915), 395--407.
\bibitem{hilbertvnn1928} D. Hilbert, J. von Neumann and L. Nordheim. \"Uber die Grundlagen der Quantenmechanik. \emph{Mathematische Annalen} 98 (1928), 1--30.
\bibitem{hilbert1930} D. Hilbert. Naturerkennen und Logik. \emph{Die Naturwissenschaften} 18 (1930), 959--963.
\bibitem{hirschman1957} I. I. Hirschman Jr. A note on entropy. \emph{American Journal of Mathematics} 79 (1957), 152--156.
\bibitem{hochreiter1997} S. Hochreiter and J. Schmidhuber. Long short term memory. \emph{Neural Computation} 9 (1997), 1735--1780.
\bibitem{hoeffding1965} W. Hoeffding. Asymptotically optimal tests for multinomial distributions. \emph{The Annals of Mathematical Statistics} 36 (1965), 369--401.
\bibitem{hoffman2014} M. D. Hoffman and A. Gelman. The No U Turn Sampler: adaptively setting path lengths in Hamiltonian Monte Carlo. \emph{Journal of Machine Learning Research} 15 (2014), 1593--1623.
\bibitem{hoffmannonnela2026} T. Hoffmann and J. P. Onnela. Unifying summary statistic selection for approximate Bayesian computation. \emph{Statistics and Computing} 36 (2026), 70.
\bibitem{holleystroock1987} R. Holley and D. Stroock. Logarithmic Sobolev inequalities and stochastic Ising models. \emph{Journal of Statistical Physics} 46 (1987), 1159--1194.
\bibitem{holt2004} C. C. Holt. Forecasting seasonals and trends by exponentially weighted moving averages. \emph{International Journal of Forecasting} 20 (2004), 5--10.
\bibitem{holway1966} L. H. Holway. New statistical models for kinetic theory: methods of construction. \emph{Physics of Fluids} 9 (1966), 1658--1673.
\bibitem{hopf1934} E. Hopf. On causality, statistics and probability. \emph{Journal of Mathematics and Physics} 13 (1934), 51--102.
\bibitem{hopf1937} E. Hopf. \emph{Ergodentheorie}. Springer, 1937.
\bibitem{hopf1950} E. Hopf. The partial differential equation $u_t+uu_x = \mu u_{xx}$. \emph{Communications on Pure and Applied Mathematics} 3 (1950), 201--230.
\bibitem{houtekamerzhang2016} P. L. Houtekamer and F. Zhang. Review of the ensemble Kalman filter for atmospheric data assimilation. \emph{Monthly Weather Review} 144 (2016), 4489--4532.
\bibitem{iba1999} Y. Iba. The Nishimori line and Bayesian statistics. \emph{Journal of Physics A: Mathematical and General} 32 (1999), 3875--3888.
\bibitem{iglesias2013} M. A. Iglesias, K. J. H. Law and A. M. Stuart. Ensemble Kalman methods for inverse problems. \emph{Inverse Problems} 29 (2013), 045001.
\bibitem{isett2018} P. Isett. A proof of Onsager's conjecture. \emph{Annals of Mathematics} 188 (2018), 871--963.
\bibitem{ising1925} E. Ising. Beitrag zur Theorie des Ferromagnetismus. \emph{Zeitschrift f\"ur Physik} 31 (1925), 253--258.
\bibitem{levin1964} B. Ja. Levin. \emph{Distribution of zeros of entire functions}. Translations of Mathematical Monographs 5, American Mathematical Society, 1964 (revised edition 1980).
\bibitem{jakobson1981} M. V. Jakobson. Absolutely continuous invariant measures for one parameter families of one dimensional maps. \emph{Communications in Mathematical Physics} 81 (1981), 39--88.
\bibitem{jamesroynetteyor2008} L. F. James, B. Roynette and M. Yor. Generalized gamma convolutions, Dirichlet means, Thorin measures, with explicit examples. \emph{Probability Surveys} 5 (2008), 346--415.
\bibitem{jamison1975} B. Jamison. The Markov processes of Schr\"odinger. \emph{Zeitschrift f\"ur Wahrscheinlichkeitstheorie und verwandte Gebiete} 32 (1975), 323--331.
\bibitem{janvresse2001} E. Janvresse. Spectral gap for Kac's model of Boltzmann equation. \emph{The Annals of Probability} 29 (2001), 288--304.
\bibitem{jarzynski1997} C. Jarzynski. Nonequilibrium equality for free energy differences. \emph{Physical Review Letters} 78 (1997), 2690--2693.
\bibitem{jarzynski1997master} C. Jarzynski. Equilibrium free energy differences from nonequilibrium measurements: a master equation approach. \emph{Physical Review E} 56 (1997), 5018--5035.
\bibitem{jaynes1957} E. T. Jaynes. Information theory and statistical mechanics. \emph{Physical Review} 106 (1957), 620--630.
\bibitem{jaynes1963} E. T. Jaynes. Information theory and statistical mechanics. In \emph{Statistical physics}, ed. K. W. Ford, pages 181--218. Benjamin, 1963.
\bibitem{jaynes2003} E. T. Jaynes. \emph{Probability theory: the logic of science}. Edited by G. L. Bretthorst, Cambridge University Press, 2003.
\bibitem{jeffreys1946} H. Jeffreys. An invariant form for the prior probability in estimation problems. \emph{Proceedings of the Royal Society of London A} 186 (1946), 453--461.
\bibitem{jerrum1992} M. Jerrum. Large cliques elude the Metropolis process. \emph{Random Structures \& Algorithms} 3 (1992), 347--359.
\bibitem{jiang2017} B. Jiang, T. Y. Wu, C. Zheng and W. H. Wong. Learning summary statistic for approximate Bayesian computation via deep neural network. \emph{Statistica Sinica} 27 (2017), 1595--1618.
\bibitem{jps2009} M. S. Johannes, N. G. Polson and J. R. Stroud. Optimal filtering of jump diffusions: extracting latent states from asset prices. \emph{Review of Financial Studies} 22 (2009), 2759--2799.
\bibitem{jko1998} R. Jordan, D. Kinderlehrer and F. Otto. The variational formulation of the Fokker Planck equation. \emph{SIAM Journal on Mathematical Analysis} 29 (1998), 1--17.
\bibitem{joycemarjoram2008} P. Joyce and P. Marjoram. Approximately sufficient statistics and Bayesian computation. \emph{Statistical Applications in Genetics and Molecular Biology} 7 (2008), 26.
\bibitem{kac1939} M. Kac. On a characterization of the normal distribution. \emph{American Journal of Mathematics} 61 (1939), 726--728.
\bibitem{kac1947} M. Kac. On the notion of recurrence in discrete stochastic processes. \emph{Bulletin of the American Mathematical Society} 53 (1947), 1002--1010.
\bibitem{kac1949} M. Kac. On distributions of certain Wiener functionals. \emph{Transactions of the American Mathematical Society} 65 (1949), 1--13.
\bibitem{kac1956} M. Kac. Foundations of kinetic theory. In \emph{Proceedings of the Third Berkeley Symposium on Mathematical Statistics and Probability}, vol. 3, pages 173--200. University of California Press, 1956.
\bibitem{kac1959} M. Kac. \emph{Probability and related topics in physical sciences}. Interscience, 1959.
\bibitem{kallenberg2002} O. Kallenberg. \emph{Foundations of modern probability}. Second edition, Springer, 2002.
\bibitem{kallenberg2005} O. Kallenberg. \emph{Probabilistic symmetries and invariance principles}. Springer, 2005.
\bibitem{kalman1960} R. E. Kalman. A new approach to linear filtering and prediction problems. \emph{Journal of Basic Engineering} 82 (1960), 35--45.
\bibitem{kalmanbucy1961} R. E. Kalman and R. S. Bucy. New results in linear filtering and prediction theory. \emph{Journal of Basic Engineering} 83 (1961), 95--108.
\bibitem{kannan1995} R. Kannan, L. Lov\'asz and M. Simonovits. Isoperimetric problems for convex bodies and a localization lemma. \emph{Discrete \& Computational Geometry} 13 (1995), 541--559.
\bibitem{kantas2015} N. Kantas, A. Doucet, S. S. Singh, J. Maciejowski and N. Chopin. On particle methods for parameter estimation in state space models. \emph{Statistical Science} 30 (2015), 328--351.
\bibitem{kantorovich1942} L. V. Kantorovich. On the translocation of masses. \emph{Doklady Akademii Nauk SSSR} 37 (1942), 199--201; English translation in \emph{Journal of Mathematical Sciences} 133 (2006), 1381--1382.
\bibitem{karatzas1991} I. Karatzas and S. E. Shreve. \emph{Brownian Motion and Stochastic Calculus}. Second edition, Graduate Texts in Mathematics 113, Springer, 1991.
\bibitem{katznelson1971} Y. Katznelson. Ergodic automorphisms of $T^n$ are Bernoulli shifts. \emph{Israel Journal of Mathematics} 10 (1971), 186--195.
\bibitem{kechris1995} A. S. Kechris. \emph{Classical Descriptive Set Theory}. Graduate Texts in Mathematics 156, Springer, 1995.
\bibitem{kelly1956} J. L. Kelly. A new interpretation of information rate. \emph{Bell System Technical Journal} 35 (1956), 917--926.
\bibitem{kennedyohagan2001} M. C. Kennedy and A. O'Hagan. Bayesian calibration of computer models. \emph{Journal of the Royal Statistical Society B} 63 (2001), 425--464.
\bibitem{kent1980} J. T. Kent. Eigenvalue expansions for diffusion hitting times. \emph{Zeitschrift f\"ur Wahrscheinlichkeitstheorie und verwandte Gebiete} 52 (1980), 309--319.
\bibitem{kestin1984} J. Kestin, K. Knierim, E. A. Mason, B. Najafi, S. T. Ro and M. Waldman. Equilibrium and transport properties of the noble gases and their mixtures at low density. \emph{Journal of Physical and Chemical Reference Data} 13 (1984), 229--303.
\bibitem{khinchin1949} A. I. Khinchin. \emph{Mathematical foundations of statistical mechanics}. Translated by G. Gamow, Dover, 1949. Russian original, 1943.
\bibitem{kipnisvaradhan1986} C. Kipnis and S. R. S. Varadhan. Central limit theorem for additive functionals of reversible Markov processes and applications to simple exclusions. \emph{Communications in Mathematical Physics} 104 (1986), 1--19.
\bibitem{kipnis1989} C. Kipnis, S. Olla and S. R. S. Varadhan. Hydrodynamics and large deviation for simple exclusion processes. \emph{Communications on Pure and Applied Mathematics} 42 (1989), 115--137.
\bibitem{kipnislandim1999} C. Kipnis and C. Landim. \emph{Scaling Limits of Interacting Particle Systems}. Springer, 1999.
\bibitem{kleijnvandervaart2006} B. J. K. Kleijn and A. W. van der Vaart. Misspecification in infinite dimensional Bayesian statistics. \emph{The Annals of Statistics} 34 (2006), 837--877.
\bibitem{klein1931} O. Klein. Zur quantenmechanischen Begr\"undung des zweiten Hauptsatzes der W\"armelehre. \emph{Zeitschrift f\"ur Physik} 72 (1931), 767--775.
\bibitem{kleinerlott2008} B. Kleiner and J. Lott. Notes on Perelman's papers. \emph{Geometry \& Topology} 12 (2008), 2587--2855.
\bibitem{knapik2013} B. T. Knapik, A. W. van der Vaart and J. H. van Zanten. Bayesian recovery of the initial condition for the heat equation. \emph{Communications in Statistics: Theory and Methods} 42 (2013), 1294--1313.
\bibitem{knothe1957} H. Knothe. Contributions to the theory of convex bodies. \emph{Michigan Mathematical Journal} 4 (1957), 39--52.
\bibitem{kolmogorov1931} A. N. Kolmogorov. \"Uber die analytischen Methoden in der Wahrscheinlichkeitsrechnung. \emph{Mathematische Annalen} 104 (1931), 415--458.
\bibitem{kolmogorov1933} A. N. Kolmogorov. \emph{Grundbegriffe der Wahrscheinlichkeitsrechnung}. Springer, 1933.
\bibitem{kolmogorov1937} A. N. Kolmogorov, I. G. Petrovsky and N. S. Piskunov. A study of the diffusion equation with increase in the amount of substance, and its application to a biological problem. In \emph{Selected Works of A. N. Kolmogorov, Volume I}, ed. V. M. Tikhomirov, pages 242--270. Kluwer, 1991. Originally in \emph{Bulletin of Moscow University, Mathematics and Mechanics} 1 (1937), 1--25.
\bibitem{kolmogorov1954} A. N. Kolmogorov. On the conservation of conditionally periodic motions under small perturbation of the Hamiltonian. \emph{Doklady Akademii Nauk SSSR} 98 (1954), 527--530.
\bibitem{kolmogorov1958} A. N. Kolmogorov. A new metric invariant of transitive dynamical systems and automorphisms of Lebesgue spaces. \emph{Doklady Akademii Nauk SSSR} 119 (1958), 861--864.
\bibitem{kolmogorov1959} A. N. Kolmogorov. Entropy per unit time as a metric invariant of automorphisms. \emph{Doklady Akademii Nauk SSSR} 124 (1959), 754--755.
\bibitem{kolmogorov1962} A. N. Kolmogorov. A refinement of previous hypotheses concerning the local structure of turbulence in a viscous incompressible fluid at high Reynolds number. \emph{Journal of Fluid Mechanics} 13 (1962), 82--85.
\bibitem{kolmogorov1957} A. N. Kolmogorov. On the representation of continuous functions of many variables by superposition of continuous functions of one variable and addition. \emph{Doklady Akademii Nauk SSSR} 114 (1957), 953--956; English translation in \emph{American Mathematical Society Translations}, Series 2, 28 (1963), 55--59.
\bibitem{kolmogorov1965} A. N. Kolmogorov. Three approaches to the quantitative definition of information. \emph{Problems of Information Transmission} 1 (1965).
\bibitem{kolmogorov1983} A. N. Kolmogorov. Combinatorial foundations of information theory and the calculus of probabilities. \emph{Russian Mathematical Surveys} 38 (1983), 29--40.
\bibitem{kolmogorov1941a} A. N. Kolmogorov. The local structure of turbulence in incompressible viscous fluid for very large Reynolds numbers. \emph{Doklady Akademii Nauk SSSR} 30 (1941), 301--305; English translation in \emph{Proceedings of the Royal Society of London A} 434 (1991), 9--13.
\bibitem{kolmogorov1941b} A. N. Kolmogorov. Dissipation of energy in the locally isotropic turbulence. \emph{Doklady Akademii Nauk SSSR} 32 (1941), 16--18; English translation in \emph{Proceedings of the Royal Society of London A} 434 (1991), 15--17.
\bibitem{konen2026} D. Konen. Inverting the Fisher information operator in nonlinear models. Preprint, arXiv:2601.13254, 2026.
\bibitem{konennickl2026} D. Konen and R. Nickl. Data assimilation with the 2D Navier Stokes equations: optimal Gaussian asymptotics for the posterior measure. \emph{The Annals of Statistics} 54 (2026), 2352--2378.
\bibitem{konennickl2026attractor} D. Konen and R. Nickl. On statistical inference for nonlinear dynamical systems evolving in their global attractor. Preprint, arXiv:2606.06018, 2026.
\bibitem{klw1994} A. Kong, J. S. Liu and W. H. Wong. Sequential imputations and Bayesian missing data problems. \emph{Journal of the American Statistical Association} 89 (1994), 278--288.
\bibitem{kontogiannis2025} A. Kontogiannis, S. V. Elgersma, A. J. Sederman and M. P. Juniper. Bayesian inverse Navier Stokes problems: joint flow field reconstruction and parameter learning. \emph{Inverse Problems} 41 (2025), 015008.
\bibitem{koopman1931} B. O. Koopman. Hamiltonian systems and transformation in Hilbert space. \emph{Proceedings of the National Academy of Sciences} 17 (1931), 315--318.
\bibitem{koopmanvonneumann1932} B. O. Koopman and J. von Neumann. Dynamical systems of continuous spectra. \emph{Proceedings of the National Academy of Sciences} 18 (1932), 255--263.
\bibitem{koopman1936} B. O. Koopman. On distributions admitting a sufficient statistic. \emph{Transactions of the American Mathematical Society} 39 (1936), 399--409.
\bibitem{koskela2020} J. Koskela, P. A. Jenkins, A. M. Johansen and D. Span\`o. Asymptotic genealogies of interacting particle systems with an application to sequential Monte Carlo. \emph{The Annals of Statistics} 48 (2020), 560--583; erratum, \emph{The Annals of Statistics} 50 (2022), 2467--2468.
\bibitem{kotani1982} S. Kotani and S. Watanabe. Krein's spectral theory of strings and generalized diffusion processes. In \emph{Functional analysis in Markov processes}, ed. M. Fukushima, pages 235--259. Lecture Notes in Mathematics 923, Springer, 1982.
\bibitem{kovachki2023} N. Kovachki, Z. Li, B. Liu, K. Azizzadenesheli, K. Bhattacharya, A. Stuart and A. Anandkumar. Neural operator: learning maps between function spaces with applications to PDEs. \emph{Journal of Machine Learning Research} 24 (2023), 1--97.
\bibitem{kraichnan1967} R. H. Kraichnan. Inertial ranges in two dimensional turbulence. \emph{Physics of Fluids} 10 (1967), 1417--1423.
\bibitem{kramers1941} H. A. Kramers and G. H. Wannier. Statistics of the two dimensional ferromagnet. Part I. \emph{Physical Review} 60 (1941), 252--262.
\bibitem{krichevsky1981} R. Krichevsky and V. K. Trofimov. The performance of universal encoding. \emph{IEEE Transactions on Information Theory} 27 (1981), 199--207.
\bibitem{kunisky2019} D. Kunisky, A. S. Wein and A. S. Bandeira. Notes on computational hardness of hypothesis testing: predictions using the low degree likelihood ratio. In \emph{Mathematical Analysis, its Applications and Computation: ISAAC 2019, Aveiro, Portugal, July 29 to August 2}, eds. P. Cerejeiras and M. Reissig, pages 1--50. Springer, 2022.
\bibitem{kurtzxiong1999} T. G. Kurtz and J. Xiong. Particle representations for a class of nonlinear SPDEs. \emph{Stochastic Processes and their Applications} 83 (1999), 103--126.
\bibitem{kurtzxiong2001} T. G. Kurtz and J. Xiong. Numerical solutions for a class of SPDEs with application to filtering. In \emph{Stochastics in Finite and Infinite Dimensions: In Honor of Gopinath Kallianpur}, eds. T. Hida, R. L. Karandikar, H. Kunita, B. S. Rajput, S. Watanabe and J. Xiong, pages 233--258. Birkh\"auser, 2001.
\bibitem{kushnerclark1978} H. J. Kushner and D. S. Clark. \emph{Stochastic Approximation Methods for Constrained and Unconstrained Systems}. Springer, 1978.
\bibitem{lamperti1972} J. Lamperti. Semistable Markov processes. I. \emph{Zeitschrift f\"ur Wahrscheinlichkeitstheorie und verwandte Gebiete} 22 (1972), 205--225.
\bibitem{landauer1961} R. Landauer. Irreversibility and heat generation in the computing process. \emph{IBM Journal of Research and Development} 5 (1961), 183--191.
\bibitem{lanford1973} O. E. Lanford. Entropy and equilibrium states in classical statistical mechanics. In \emph{Statistical Mechanics and Mathematical Problems}, ed. A. Lenard, Lecture Notes in Physics 20, pages 1--113. Springer, 1973.
\bibitem{lanford1975} O. E. Lanford. Time evolution of large classical systems. In \emph{Dynamical systems, theory and applications}, Lecture Notes in Physics 38, pages 1--111. Springer, 1975.
\bibitem{langevin1908} P. Langevin. Sur la th\'eorie du mouvement brownien. \emph{Comptes Rendus de l'Acad\'emie des Sciences} 146 (1908), 530--533.
\bibitem{laplace1814} P. S. Laplace. \emph{Essai philosophique sur les probabilit\'es}. Courcier, Paris, 1814.
\bibitem{lauritzen1981} S. L. Lauritzen. Time series analysis in 1880: a discussion of contributions made by T. N. Thiele. \emph{International Statistical Review} 49 (1981), 319--331.
\bibitem{law2015} K. Law, A. Stuart and K. Zygalakis. \emph{Data assimilation: a mathematical introduction}. Springer, 2015.
\bibitem{lebesgue1902} H. Lebesgue. Int\'egrale, longueur, aire. \emph{Annali di Matematica Pura ed Applicata} 7 (1902), 231--359.
\bibitem{ledimettalagrand1986} F. X. Le Dimet and O. Talagrand. Variational algorithms for analysis and assimilation of meteorological observations: theoretical aspects. \emph{Tellus A} 38 (1986), 97--110.
\bibitem{ledoux2001} M. Ledoux. \emph{The Concentration of Measure Phenomenon}. Mathematical Surveys and Monographs 89, American Mathematical Society, 2001.
\bibitem{ledrappier1981} F. Ledrappier. Some properties of absolutely continuous invariant measures on an interval. \emph{Ergodic Theory and Dynamical Systems} 1 (1981), 77--93.
\bibitem{leeyang1952} T. D. Lee and C. N. Yang. Statistical theory of equations of state and phase transitions. II. Lattice gas and Ising model. \emph{Physical Review} 87 (1952), 410--419.
\bibitem{lee2002} H. K. H. Lee, D. M. Higdon, Z. Bi, M. A. R. Ferreira and M. West. Markov random field models for high dimensional parameters in simulations of fluid flow in porous media. \emph{Technometrics} 44 (2002), 230--241.
\bibitem{leffrex1990} H. S. Leff and A. F. Rex, eds. \emph{Maxwell's Demon: Entropy, Information, Computing}. Princeton University Press and Adam Hilger, 1990.
\bibitem{legland2011} F. Le Gland, V. Monbet and V. D. Tran. Large sample asymptotics for the ensemble Kalman filter. In \emph{The Oxford Handbook of Nonlinear Filtering}, eds. D. Crisan and B. Rozovskii, pages 598--631. Oxford University Press, 2011.
\bibitem{lehec2013} J. Lehec. Representation formula for the entropy and functional inequalities. \emph{Annales de l'Institut Henri Poincar\'e, Probabilit\'es et Statistiques} 49 (2013), 885--899.
\bibitem{leith1968} C. E. Leith. Diffusion approximation for two dimensional turbulence. \emph{Physics of Fluids} 11 (1968), 671--672.
\bibitem{lelarge2019} M. Lelarge and L. Miolane. Fundamental limits of symmetric low rank matrix estimation. \emph{Probability Theory and Related Fields} 173 (2019), 859--929.
\bibitem{leonard2014} C. L\'eonard. A survey of the Schr\"odinger problem and some of its connections with optimal transport. \emph{Discrete and Continuous Dynamical Systems A} 34 (2014), 1533--1574.
\bibitem{leonard2012} C. L\'eonard. From the Schr\"odinger problem to the Monge Kantorovich problem. \emph{Journal of Functional Analysis} 262 (2012), 1879--1920.
\bibitem{leonard2010} C. L\'eonard. Entropic projections and dominating points. \emph{ESAIM: Probability and Statistics} 14 (2010), 343--381.
\bibitem{leray1934} J. Leray. Sur le mouvement d'un liquide visqueux emplissant l'espace. \emph{Acta Mathematica} 63 (1934), 193--248.
\bibitem{lesieur2017} T. Lesieur, F. Krzakala and L. Zdeborov\'a. Constrained low rank matrix estimation: phase transitions, approximate message passing and applications. \emph{Journal of Statistical Mechanics: Theory and Experiment} (2017), 073403.
\bibitem{levermore1996} C. D. Levermore. Moment closure hierarchies for kinetic theories. \emph{Journal of Statistical Physics} 83 (1996), 1021--1065.
\bibitem{levermoremasmoudi2010} C. D. Levermore and N. Masmoudi. From the Boltzmann equation to an incompressible Navier Stokes Fourier system. \emph{Archive for Rational Mechanics and Analysis} 196 (2010), 753--809.
\bibitem{levin1973} L. A. Levin. On the notion of a random sequence. \emph{Soviet Mathematics Doklady} 14 (1973), 1413--1416.
\bibitem{levin1974} L. A. Levin. Laws of information conservation (nongrowth) and aspects of the foundation of probability theory. \emph{Problems of Information Transmission} 10 (1974), 206--210.
\bibitem{levycalvert2021} W. B. Levy and V. G. Calvert. Communication consumes 35 times more energy than computation in the human cortex, but both costs are needed to predict synapse number. \emph{Proceedings of the National Academy of Sciences} 118 (2021), e2008173118.
\bibitem{liduan1989} K. C. Li and N. Duan. Regression analysis under link violation. \emph{The Annals of Statistics} 17 (1989), 1009--1052.
\bibitem{livitanyi2008} M. Li and P. Vit\'anyi. \emph{An Introduction to Kolmogorov Complexity and Its Applications}. Third edition, Springer, 2008.
\bibitem{li2021} Z. Li, N. Kovachki, K. Azizzadenesheli, B. Liu, K. Bhattacharya, A. Stuart and A. Anandkumar. Fourier neural operator for parametric partial differential equations. In \emph{International Conference on Learning Representations}, 2021.
\bibitem{liggett1985} T. M. Liggett. \emph{Interacting particle systems}. Springer, 1985.
\bibitem{lindblad1975} G. Lindblad. Completely positive maps and entropy inequalities. \emph{Communications in Mathematical Physics} 40 (1975), 147--151.
\bibitem{lindgrenrue2011} F. Lindgren, H. Rue and J. Lindstr\"om. An explicit link between Gaussian fields and Gaussian Markov random fields: the stochastic partial differential equation approach. \emph{Journal of the Royal Statistical Society B} 73 (2011), 423--498.
\bibitem{lindley1956} D. V. Lindley. On a measure of the information provided by an experiment. \emph{The Annals of Mathematical Statistics} 27 (1956), 986--1005.
\bibitem{lions1969} J. L. Lions. \emph{Quelques m\'ethodes de r\'esolution des probl\`emes aux limites non lin\'eaires}. Dunod, 1969.
\bibitem{lipman2023} Y. Lipman, R. T. Q. Chen, H. Ben Hamu, M. Nickel and M. Le. Flow matching for generative modeling. \emph{International Conference on Learning Representations} (2023).
\bibitem{litowitz2026} A. Litowitz, N. G. Polson and V. Sokolov. Knowledge, LLMs, and the information ladder: from Shannon to Kolmogorov. Working paper, SSRN 6595298, 2026.
\bibitem{liuwest2001} J. Liu and M. West. Combined parameter and state estimation in simulation based filtering. In \emph{Sequential Monte Carlo methods in practice}, eds. A. Doucet, N. de Freitas and N. Gordon, pages 197--223. Springer, 2001.
\bibitem{liu2024kan} Z. Liu, Y. Wang, S. Vaidya, F. Ruehle, J. Halverson, M. Solja\v{c}i\'c, T. Y. Hou and M. Tegmark. KAN: Kolmogorov Arnold networks. \emph{International Conference on Learning Representations} (2025).
\bibitem{ljung1977} L. Ljung. Analysis of recursive stochastic algorithms. \emph{IEEE Transactions on Automatic Control} 22 (1977), 551--575.
\bibitem{longstaffschwartz2001} F. A. Longstaff and E. S. Schwartz. Valuing American options by simulation: a simple least squares approach. \emph{Review of Financial Studies} 14 (2001), 113--147.
\bibitem{weightedbootstrap2026} H. F. Lopes, N. G. Polson and V. Sokolov. Uncertainty quantification: from weighted bootstrap to generators. Working paper, 2026.
\bibitem{loschmidt1876} J. Loschmidt. \"Uber den Zustand des W\"armegleichgewichtes eines Systems von K\"orpern mit R\"ucksicht auf die Schwerkraft. \emph{Sitzungsberichte der Akademie der Wissenschaften Wien} 73 (1876), 128--142.
\bibitem{lottvillani2009} J. Lott and C. Villani. Ricci curvature for metric measure spaces via optimal transport. \emph{Annals of Mathematics} 169 (2009), 903--991.
\bibitem{lovasz1993} L. Lov\'asz and M. Simonovits. Random walks in a convex body and an improved volume algorithm. \emph{Random Structures and Algorithms} 4 (1993), 359--412.
\bibitem{luders1950} G. L\"uders. \"Uber die Zustands\"anderung durch den Me\ss{}proze\ss. \emph{Annalen der Physik} 8 (1951), 322--328.
\bibitem{luohou2014} G. Luo and T. Y. Hou. Potentially singular solutions of the 3D axisymmetric Euler equations. \emph{Proceedings of the National Academy of Sciences} 111 (2014), 12968--12973.
\bibitem{lyubich2002} M. Lyubich. Almost every real quadratic map is either regular or stochastic. \emph{Annals of Mathematics} 156 (2002), 1--78.
\bibitem{ma2019} Y. A. Ma, Y. Chen, C. Jin, N. Flammarion and M. I. Jordan. Sampling can be faster than optimization. \emph{Proceedings of the National Academy of Sciences} 116 (2019), 20881--20885.
\bibitem{mackey1963} G. W. Mackey. \emph{Mathematical foundations of quantum mechanics: a lecture note volume}. W. A. Benjamin, 1963.
\bibitem{madelung1927} E. Madelung. Quantentheorie in hydrodynamischer Form. \emph{Zeitschrift f\"ur Physik} 40 (1927), 322--326.
\bibitem{mandt2017} S. Mandt, M. D. Hoffman and D. M. Blei. Stochastic gradient descent as approximate Bayesian inference. \emph{Journal of Machine Learning Research} 18 (2017), 1--35.
\bibitem{mansuymazliak2005} R. Mansuy and L. Mazliak. Introduction au rapport de Poincar\'e pour le proc\`es en cassation de Dreyfus en 1904. \emph{Journal \'Electronique d'Histoire des Probabilit\'es et de la Statistique} 1 (2005), 1--4.
\bibitem{marklof2011} J. Marklof and A. Str\"ombergsson. The Boltzmann Grad limit of the periodic Lorentz gas. \emph{Annals of Mathematics} 174 (2011), 225--298.
\bibitem{markusmeyer1974} L. Markus and K. R. Meyer. Generic Hamiltonian dynamical systems are neither integrable nor ergodic. \emph{Memoirs of the American Mathematical Society} 144 (1974).
\bibitem{martinlof1966} P. Martin L\"of. The definition of random sequences. \emph{Information and Control} 9 (1966), 602--619.
\bibitem{martinlof1971} P. Martin L\"of. Complexity oscillations in infinite binary sequences. \emph{Zeitschrift f\"ur Wahrscheinlichkeitstheorie und verwandte Gebiete} 19 (1971), 225--230.
\bibitem{masreliez1975} C. J. Masreliez. Approximate non Gaussian filtering with linear state and observation relations. \emph{IEEE Transactions on Automatic Control} 20 (1975), 107--110.
\bibitem{matsumotoyor2005} H. Matsumoto and M. Yor. Exponential functionals of Brownian motion I and II. \emph{Probability Surveys} 2 (2005), 312--384.
\bibitem{maxwell1860} J. C. Maxwell. Illustrations of the dynamical theory of gases. Part I. On the motions and collisions of perfectly elastic spheres. \emph{Philosophical Magazine} 19 (1860), 19--32.
\bibitem{maxwell1867} J. C. Maxwell. On the dynamical theory of gases. \emph{Philosophical Transactions of the Royal Society of London} 157 (1867), 49--88.
\bibitem{maxwell1871} J. C. Maxwell. \emph{Theory of Heat}. Longmans, Green and Co., London, 1871.
\bibitem{mazliak2012} L. Mazliak. Poincar\'e's odds. In \emph{Henri Poincar\'e, 1912 to 2012: Poincar\'e Seminar 2012}, pages 151--192. Springer Basel, 2015.
\bibitem{mazurvanderlinden1963} P. Mazur and J. van der Linden. Asymptotic form of the structure function for real systems. \emph{Journal of Mathematical Physics} 4 (1963), 271--277.
\bibitem{mccann1995} R. J. McCann. Existence and uniqueness of monotone measure preserving maps. \emph{Duke Mathematical Journal} 80 (1995), 309--323.
\bibitem{mccann1997} R. J. McCann. A convexity principle for interacting gases. \emph{Advances in Mathematics} 128 (1997), 153--179.
\bibitem{mccormmach1967} R. McCormmach. Henri Poincar\'e and the quantum theory. \emph{Isis} 58 (1967), 37--55.
\bibitem{mckean1966} H. P. McKean Jr. Speed of approach to equilibrium for Kac's caricature of a Maxwellian gas. \emph{Archive for Rational Mechanics and Analysis} 21 (1966), 343--367.
\bibitem{mcmillan1953} B. McMillan. The basic theorems of information theory. \emph{The Annals of Mathematical Statistics} 24 (1953), 196--219.
\bibitem{menaniles2019} G. Mena and J. Niles Weed. Statistical bounds for entropic optimal transport: sample complexity and the central limit theorem. \emph{Advances in Neural Information Processing Systems} 32 (2019).
\bibitem{metropolis1953} N. Metropolis, A. W. Rosenbluth, M. N. Rosenbluth, A. H. Teller and E. Teller. Equation of state calculations by fast computing machines. \emph{Journal of Chemical Physics} 21 (1953), 1087--1092.
\bibitem{mezard2009} M. M\'ezard and A. Montanari. \emph{Information, Physics, and Computation}. Oxford University Press, 2009.
\bibitem{michael1976} J. R. Michael, W. R. Schucany and R. W. Haas. Generating random variates using transformations with multiple roots. \emph{The American Statistician} 30 (1976), 88--90.
\bibitem{mikami2004} T. Mikami. Monge's problem with a quadratic cost by the zero noise limit of h path processes. \emph{Probability Theory and Related Fields} 129 (2004), 245--260.
\bibitem{mischlermouhot2013} S. Mischler and C. Mouhot. Kac's program in kinetic theory. \emph{Inventiones Mathematicae} 193 (2013), 1--147.
\bibitem{morgan2007} J. Morgan and G. Tian. \emph{Ricci Flow and the Poincar\'e Conjecture}. American Mathematical Society and Clay Mathematics Institute, 2007.
\bibitem{moser1960} J. Moser. A new proof of De Giorgi's theorem concerning the regularity problem for elliptic differential equations. \emph{Communications on Pure and Applied Mathematics} 13 (1960), 457--468.
\bibitem{moser1962} J. Moser. On invariant curves of area preserving mappings of an annulus. \emph{Nachrichten der Akademie der Wissenschaften in G\"ottingen, Mathematisch Physikalische Klasse} II (1962), 1--20.
\bibitem{mouhotvillani2011} C. Mouhot and C. Villani. On Landau damping. \emph{Acta Mathematica} 207 (2011), 29--201.
\bibitem{muth1960} J. F. Muth. Optimal properties of exponentially weighted forecasts. \emph{Journal of the American Statistical Association} 55 (1960), 299--306.
\bibitem{nareklishvili2023} M. Nareklishvili, N. G. Polson and V. Sokolov. Deep partial least squares for instrumental variable regression. \emph{Applied Stochastic Models in Business and Industry} 39 (2023), 734--754.
\bibitem{nareklishvili2025review} M. Nareklishvili, N. G. Polson and V. Sokolov. Generative modeling: a review. Preprint, arXiv:2501.05458, 2025.
\bibitem{nash1958} J. Nash. Continuity of solutions of parabolic and elliptic equations. \emph{American Journal of Mathematics} 80 (1958), 931--954.
\bibitem{neal2001} R. M. Neal. Annealed importance sampling. \emph{Statistics and Computing} 11 (2001), 125--139.
\bibitem{neal2003} R. M. Neal. Slice sampling. \emph{The Annals of Statistics} 31 (2003), 705--767.
\bibitem{neal2011} R. M. Neal. MCMC using Hamiltonian dynamics. In \emph{Handbook of Markov Chain Monte Carlo}, eds. S. Brooks, A. Gelman, G. Jones and X. L. Meng, pages 113--162. Chapman and Hall/CRC, 2011.
\bibitem{necas1996} J. Ne\v{c}as, M. R\r{u}\v{z}i\v{c}ka and V. \v{S}ver\'ak. On Leray's self similar solutions of the Navier Stokes equations. \emph{Acta Mathematica} 176 (1996), 283--294.
\bibitem{nelson1966} E. Nelson. Derivation of the Schr\"odinger equation from Newtonian mechanics. \emph{Physical Review} 150 (1966), 1079--1085.
\bibitem{nelson1967} E. Nelson. \emph{Dynamical theories of Brownian motion}. Princeton University Press, 1967.
\bibitem{nelson1973} E. Nelson. The free Markoff field. \emph{Journal of Functional Analysis} 12 (1973), 211--227.
\bibitem{nelson1985} E. Nelson. \emph{Quantum Fluctuations}. Princeton University Press, 1985.
\bibitem{nelson1992} D. B. Nelson. Filtering and forecasting with misspecified ARCH models I: getting the right variance with the wrong model. \emph{Journal of Econometrics} 52 (1992), 61--90.
\bibitem{nelson2012} E. Nelson. Review of stochastic mechanics. \emph{Journal of Physics: Conference Series} 361 (2012), 012011.
\bibitem{newman1974} C. M. Newman. Zeros of the partition function for generalized Ising systems. \emph{Communications on Pure and Applied Mathematics} 27 (1974), 143--159.
\bibitem{newman1975} C. M. Newman. Inequalities for Ising models and field theories which obey the Lee Yang theorem. \emph{Communications in Mathematical Physics} 41 (1975), 1--9.
\bibitem{newman1976} C. M. Newman. Fourier transforms with only real zeros. \emph{Proceedings of the American Mathematical Society} 61 (1976), 245--251.
\bibitem{newtonraftery1994} M. A. Newton and A. E. Raftery. Approximate Bayesian inference with the weighted likelihood bootstrap. \emph{Journal of the Royal Statistical Society B} 56 (1994), 3--26.
\bibitem{ni2004} L. Ni. The entropy formula for linear heat equation. \emph{Journal of Geometric Analysis} 14 (2004), 87--100.
\bibitem{nickltiti2024} R. Nickl and E. S. Titi. On posterior consistency of data assimilation with Gaussian process priors: the 2D Navier Stokes equations. \emph{The Annals of Statistics} 52 (2024), 1825--1844.
\bibitem{nickl2023book} R. Nickl. \emph{Bayesian nonlinear statistical inverse problems}. Zurich Lectures in Advanced Mathematics, European Mathematical Society, 2023.
\bibitem{nicklbvm2024} R. Nickl. Bernstein von Mises theorems for time evolution equations. Preprint, arXiv:2407.14781, 2024.
\bibitem{nicklwang2024} R. Nickl and S. Wang. On polynomial time computation of high dimensional posterior measures by Langevin type algorithms. \emph{Journal of the European Mathematical Society} 26 (2024), 1031--1112.
\bibitem{nickl2025} R. Nickl, G. A. Pavliotis and K. Ray. Bayesian nonparametric inference in McKean Vlasov models. \emph{The Annals of Statistics} 53 (2025), 170--193.
\bibitem{nishimori1981} H. Nishimori. Internal energy, specific heat and correlation function of the bond random Ising model. \emph{Progress of Theoretical Physics} 66 (1981), 1169--1181.
\bibitem{nunesbalding2010} M. A. Nunes and D. J. Balding. On optimal selection of summary statistics for approximate Bayesian computation. \emph{Statistical Applications in Genetics and Molecular Biology} 9 (2010), 34.
\bibitem{obukhov1962} A. M. Obukhov. Some specific features of atmospheric turbulence. \emph{Journal of Fluid Mechanics} 13 (1962), 77--81.
\bibitem{oja1982} E. Oja. Simplified neuron model as a principal component analyzer. \emph{Journal of Mathematical Biology} 15 (1982), 267--273.
\bibitem{dlmf} F. W. J. Olver, A. B. Olde Daalhuis, D. W. Lozier, B. I. Schneider, R. F. Boisvert, C. W. Clark, B. R. Miller, B. V. Saunders, H. S. Cohl and M. A. McClain, eds. \emph{NIST Digital Library of Mathematical Functions}. National Institute of Standards and Technology, Release 1.2.8 of 15 September 2026, dlmf.nist.gov.
\bibitem{onsager1944} L. Onsager. Crystal statistics. I. A two dimensional model with an order disorder transition. \emph{Physical Review} 65 (1944), 117--149.
\bibitem{onsager1949} L. Onsager. Statistical hydrodynamics. \emph{Il Nuovo Cimento} 6, Supplement 2 (1949), 279--287.
\bibitem{ornstein1970} D. S. Ornstein. Bernoulli shifts with the same entropy are isomorphic. \emph{Advances in Mathematics} 4 (1970), 337--352.
\bibitem{ornsteinweiss1973} D. Ornstein and B. Weiss. Geodesic flows are Bernoullian. \emph{Israel Journal of Mathematics} 14 (1973), 184--198.
\bibitem{osterwalder1973} K. Osterwalder and R. Schrader. Axioms for Euclidean Green's functions. \emph{Communications in Mathematical Physics} 31 (1973), 83--112.
\bibitem{osterwalder1975} K. Osterwalder and R. Schrader. Axioms for Euclidean Green's functions II. \emph{Communications in Mathematical Physics} 42 (1975), 281--305.
\bibitem{ottovillani2000} F. Otto and C. Villani. Generalization of an inequality by Talagrand and links with the logarithmic Sobolev inequality. \emph{Journal of Functional Analysis} 173 (2000), 361--400.
\bibitem{otto2001} F. Otto. The geometry of dissipative evolution equations: the porous medium equation. \emph{Communications in Partial Differential Equations} 26 (2001), 101--174.
\bibitem{pacchiardi2024} L. Pacchiardi, S. Khoo and R. Dutta. Generalized Bayesian likelihood free inference. \emph{Electronic Journal of Statistics} 18 (2024), 3628--3686.
\bibitem{padilla2022} O. H. M. Padilla, W. Tansey and Y. Chen. Quantile regression with ReLU networks: estimators and minimax rates. \emph{Journal of Machine Learning Research} 23 (2022), 1--42.
\bibitem{papamakarios2021} G. Papamakarios, E. Nalisnick, D. J. Rezende, S. Mohamed and B. Lakshminarayanan. Normalizing flows for probabilistic modeling and inference. \emph{Journal of Machine Learning Research} 22 (2021), 1--64.
\bibitem{park2016} M. Park, W. Jitkrittum and D. Sejdinovic. K2 ABC: approximate Bayesian computation with kernel embeddings. In \emph{Proceedings of the 19th International Conference on Artificial Intelligence and Statistics}, PMLR 51, pages 398--407, 2016.
\bibitem{pearl1982} J. Pearl. Reverend Bayes on inference engines: a distributed hierarchical approach. In \emph{Proceedings of the National Conference on Artificial Intelligence (AAAI 82)}, pages 133--136, 1982.
\bibitem{pedersen1995} A. R. Pedersen. A new approach to maximum likelihood estimation for stochastic differential equations based on discrete observations. \emph{Scandinavian Journal of Statistics} 22 (1995), 55--71.
\bibitem{peierls1936} R. Peierls. On Ising's model of ferromagnetism. \emph{Mathematical Proceedings of the Cambridge Philosophical Society} 32 (1936), 477--481.
\bibitem{perelman2002} G. Perelman. The entropy formula for the Ricci flow and its geometric applications. Preprint, arXiv:math/0211159, 2002.
\bibitem{perelman2003a} G. Perelman. Ricci flow with surgery on three manifolds. Preprint, arXiv:math/0303109, 2003.
\bibitem{perelman2003b} G. Perelman. Finite extinction time for the solutions to the Ricci flow on certain three manifolds. \emph{Topologica} 1 (2008), 005.
\bibitem{perrin1908} J. Perrin. L'agitation mol\'eculaire et le mouvement brownien. \emph{Comptes Rendus de l'Acad\'emie des Sciences} 146 (1908), 967--970.
\bibitem{perrin1909} J. Perrin. Mouvement brownien et r\'ealit\'e mol\'eculaire. \emph{Annales de Chimie et de Physique}, 8th series, 18 (1909), 5--114.
\bibitem{perrin1926} J. B. Perrin. Discontinuous structure of matter. Nobel lecture, 11 December 1926. In \emph{Nobel Lectures, Physics 1922 to 1941}, pages 138--164. Elsevier, 1965.
\bibitem{petersen1983} K. Petersen. \emph{Ergodic theory}. Cambridge Studies in Advanced Mathematics 2, Cambridge University Press, 1983.
\bibitem{petz1986} D. Petz. Sufficient subalgebras and the relative entropy of states of a von Neumann algebra. \emph{Communications in Mathematical Physics} 105 (1986), 123--131.
\bibitem{petz1988} D. Petz. Sufficiency of channels over von Neumann algebras. \emph{Quarterly Journal of Mathematics} 39 (1988), 97--108.
\bibitem{peyrecuturi2019} G. Peyr\'e and M. Cuturi. Computational optimal transport with applications to data sciences. \emph{Foundations and Trends in Machine Learning} 11 (2019), 355--607.
\bibitem{phelps2001} R. R. Phelps. \emph{Lectures on Choquet's theorem}. Second edition, Lecture Notes in Mathematics 1757, Springer, 2001.
\bibitem{pitman1936} E. J. G. Pitman. Sufficient statistics and intrinsic accuracy. \emph{Mathematical Proceedings of the Cambridge Philosophical Society} 32 (1936), 567--579.
\bibitem{pitman1975} J. W. Pitman. One dimensional Brownian motion and the three dimensional Bessel process. \emph{Advances in Applied Probability} 7 (1975), 511--526.
\bibitem{pitman1997} J. Pitman and M. Yor. The two parameter Poisson Dirichlet distribution derived from a stable subordinator. \emph{The Annals of Probability} 25 (1997), 855--900.
\bibitem{pitman2003} J. Pitman and M. Yor. Infinitely divisible laws associated with hyperbolic functions. \emph{Canadian Journal of Mathematics} 55 (2003), 292--330.
\bibitem{pittshephard1999} M. K. Pitt and N. Shephard. Filtering via simulation: auxiliary particle filters. \emph{Journal of the American Statistical Association} 94 (1999), 590--599.
\bibitem{pitt2012} M. K. Pitt, R. dos Santos Silva, P. Giordani and R. Kohn. On some properties of Markov chain Monte Carlo simulation methods based on the particle filter. \emph{Journal of Econometrics} 171 (2012), 134--151.
\bibitem{plancherel1913} M. Plancherel. Beweis der Unm\"oglichkeit ergodischer mechanischer Systeme. \emph{Annalen der Physik} 42 (1913), 1061--1063.
\bibitem{ptv2017} F. Planchon, N. Tzvetkov and N. Visciglia. On the growth of Sobolev norms for NLS on 2 and 3 dimensional manifolds. \emph{Analysis \& PDE} 10 (2017), 1123--1147.
\bibitem{planck1901} M. Planck. Ueber das Gesetz der Energieverteilung im Normalspectrum. \emph{Annalen der Physik} 4 (1901), 553--563.
\bibitem{platttrudgian2021} D. Platt and T. Trudgian. The Riemann hypothesis is true up to $3\cdot10^{12}$. \emph{Bulletin of the London Mathematical Society} 53 (2021), 792--797.
\bibitem{poincare1890} H. Poincar\'e. Sur le probl\`eme des trois corps et les \'equations de la dynamique. \emph{Acta Mathematica} 13 (1890), 1--270.
\bibitem{poincare1890b} H. Poincar\'e. Sur les \'equations aux d\'eriv\'ees partielles de la physique math\'ematique. \emph{American Journal of Mathematics} 12 (1890), 211--294.
\bibitem{poincare1893} H. Poincar\'e. Le m\'ecanisme et l'exp\'erience. \emph{Revue de M\'etaphysique et de Morale} 1 (1893), 534--537.
\bibitem{poincare1899} H. Poincar\'e. R\'eflexions sur le calcul des probabilit\'es. \emph{Revue g\'en\'erale des sciences pures et appliqu\'ees} 10 (1899), 262--269.
\bibitem{poincare1902} H. Poincar\'e. \emph{La science et l'hypoth\`ese}. Flammarion, 1902.
\bibitem{poincare1904} H. Poincar\'e. Cinqui\`eme compl\'ement \`a l'Analysis situs. \emph{Rendiconti del Circolo Matematico di Palermo} 18 (1904), 45--110.
\bibitem{poincare1906} H. Poincar\'e. R\'eflexions sur la th\'eorie cin\'etique des gaz. \emph{Journal de Physique Th\'eorique et Appliqu\'ee} 5 (1906), 369--403.
\bibitem{poincare1912calcul} H. Poincar\'e. \emph{Calcul des probabilit\'es}. Second edition, Gauthier Villars, 1912.
\bibitem{poincare1912quanta} H. Poincar\'e. Sur la th\'eorie des quanta. \emph{Journal de Physique Th\'eorique et Appliqu\'ee} 2 (1912), 5--34.
\bibitem{polson1992} N. G. Polson. Comment on ``Practical Markov chain Monte Carlo'' by C. J. Geyer. \emph{Statistical Science} 7 (1992), 490--491.
\bibitem{polsonroberts1993} N. G. Polson and G. O. Roberts. A utility based approach to information for stochastic differential equations. \emph{Stochastic Processes and their Applications} 48 (1993), 341--356.
\bibitem{polsonroberts1994} N. G. Polson and G. O. Roberts. Bayes factors for discrete observations from diffusion processes. \emph{Biometrika} 81 (1994), 11--26.
\bibitem{psm2008} N. G. Polson, J. R. Stroud and P. M\"uller. Practical filtering with sequential parameter learning. \emph{Journal of the Royal Statistical Society B} 70 (2008), 413--428.
\bibitem{psw2013} N. G. Polson, J. G. Scott and J. Windle. Bayesian inference for logistic models using P\'olya Gamma latent variables. \emph{Journal of the American Statistical Association} 108 (2013), 1339--1349.
\bibitem{polsonrockova2018} N. G. Polson and V. Ro\v{c}kov\'a. Posterior concentration for sparse deep learning. In \emph{Advances in Neural Information Processing Systems 31}, pages 930--941, 2018.
\bibitem{polsonsokolovxu2021} N. G. Polson, V. Sokolov and J. Xu. Deep learning partial least squares. Preprint, arXiv:2106.14085, 2021.
\bibitem{polsonruggerisokolov2024} N. G. Polson, F. Ruggeri and V. Sokolov. Generative Bayesian computation for maximum expected utility. \emph{Entropy} 26 (2024), 1076.
\bibitem{polsonsokolov2025} N. G. Polson and V. Sokolov. Generative AI for Bayesian computation. \emph{Entropy} 27 (2025), 683.
\bibitem{polsonsokolov2025dd} N. G. Polson and V. Sokolov. Bayesian double descent. Preprint, arXiv:2507.07338, 2025.
\bibitem{polsonh8} N. G. Polson. On Hilbert's 8th problem. Preprint, arXiv:1708.02653, 2017 (revised September 2026).
\bibitem{polsonsokolov2026} N. G. Polson and V. Sokolov. Can generative Bayesian computation beat the Bernstein von Mises theorem? A numerical and theoretical study of the Konen Nickl two dimensional Navier Stokes problem. Companion numerical study, unpublished, 2026.
\bibitem{polsonsokolov2026air} N. G. Polson and V. Sokolov. AI assisted research is AlphaZero: digital twins, self play and recursive self learning. Working paper, 2026.
\bibitem{polsonsokolov2026b} N. G. Polson and V. Sokolov. \emph{Bayes, AI, and Deep Learning: Foundations of Data Science}. Chapman \& Hall/CRC, 2026.
\bibitem{polsonsokolov2026gates} N. G. Polson and V. Sokolov. Gates are Bayes: neural architectures as approximate nonlinear filters. Preprint, SSRN 7254828, 2026.
\bibitem{polsonsokolov2026gbcshot} N. G. Polson and V. Sokolov. GBC for shot noise processes. Preprint, SSRN 7035698, 2026.
\bibitem{polsonsokolov2026manifold} N. G. Polson and V. Sokolov. Bayesian estimation on the solution manifold of a PDE: generative computation and adjoint sufficient statistics. Working paper, 2026.
\bibitem{polsonsokolov2026ns} N. G. Polson and V. Sokolov. Bayesian methods for the Navier Stokes equations. Preprint, arXiv:2602.02945, 2026.
\bibitem{polsonsokolov2026twentytwobits} N. G. Polson and V. Sokolov. Twenty Two Bits: entropy, fractional Kelly and the stochastic growth of a weighted coin. Working paper, 2026.
\bibitem{polsonsokolovelements} N. G. Polson and V. Sokolov. \emph{Elements of Bayes and Information Theory}. 2026.
\bibitem{polsonsokolovperelman} N. G. Polson, V. Sokolov and D. Zantedeschi. Perelman's entropy, dual densities, and the Heisenberg uncertainty principle. Working paper, SSRN 6903859, 2026.
\bibitem{polsonsokolovyor} N. G. Polson and V. Sokolov. Yor's identities as algorithms. SSRN 6789381, 2026.
\bibitem{polya1926} G. P\'olya. Bemerkung \"uber die Integraldarstellung der Riemannschen $\xi$ Funktion. \emph{Acta Mathematica} 48 (1926), 305--317.
\bibitem{polya1927} G. P\'olya. \"Uber trigonometrische Integrale mit nur reellen Nullstellen. \emph{Journal f\"ur die reine und angewandte Mathematik} 158 (1927), 6--18.
\bibitem{polymath2019} D. H. J. Polymath. Effective approximation of heat flow evolution of the Riemann $\xi$ function, and a new upper bound for the de Bruijn Newman constant. \emph{Research in the Mathematical Sciences} 6 (2019), 31.
\bibitem{prangle2018} D. Prangle. Summary statistics. In \emph{Handbook of approximate Bayesian computation}, pages 125--152. Chapman \& Hall/CRC, 2018.
\bibitem{prentis1995} J. J. Prentis. Poincar\'e's proof of the quantum discontinuity of nature. \emph{American Journal of Physics} 63 (1995), 339--350.
\bibitem{pritchard1999} J. K. Pritchard, M. T. Seielstad, A. Perez Lezaun and M. W. Feldman. Population growth of human Y chromosomes: a study of Y chromosome microsatellites. \emph{Molecular Biology and Evolution} 16 (1999), 1791--1798.
\bibitem{raissi2017} M. Raissi, P. Perdikaris and G. E. Karniadakis. Machine learning of linear differential equations using Gaussian processes. \emph{Journal of Computational Physics} 348 (2017), 683--693.
\bibitem{raynal2019} L. Raynal, J. M. Marin, P. Pudlo, M. Ribatet, C. P. Robert and A. Estoup. ABC random forests for Bayesian parameter inference. \emph{Bioinformatics} 35 (2019), 1720--1728.
\bibitem{raynal2022} L. Raynal, S. Chen, A. Mira and J. P. Onnela. Scalable approximate Bayesian computation for growing network models via extrapolated and sampled summaries. \emph{Bayesian Analysis} 17 (2022), 165--192.
\bibitem{regazzini2003} E. Regazzini, A. Lijoi and I. Pr\"unster. Distributional results for means of normalized random measures with independent increments. \emph{The Annals of Statistics} 31 (2003), 560--585.
\bibitem{revuzyor} D. Revuz and M. Yor. \emph{Continuous martingales and Brownian motion}. Third edition, Springer, 1999.
\bibitem{rice1944} S. O. Rice. Mathematical analysis of random noise. \emph{Bell System Technical Journal} 23 (1944), 282--332.
\bibitem{richardson1922} L. F. Richardson. \emph{Weather prediction by numerical process}. Cambridge University Press, 1922.
\bibitem{riemann1859} B. Riemann. \"Uber die Anzahl der Primzahlen unter einer gegebenen Gr\"osse. \emph{Monatsberichte der K\"oniglich Preussischen Akademie der Wissenschaften zu Berlin} (1859), 671--680.
\bibitem{rissanen1978} J. Rissanen. Modeling by shortest data description. \emph{Automatica} 14 (1978), 465--471.
\bibitem{rissanen1984} J. Rissanen. Universal coding, information, prediction, and estimation. \emph{IEEE Transactions on Information Theory} 30 (1984), 629--636.
\bibitem{rissanen1996} J. J. Rissanen. Fisher information and stochastic complexity. \emph{IEEE Transactions on Information Theory} 42 (1996), 40--47.
\bibitem{robbinsmonro1951} H. Robbins and S. Monro. A stochastic approximation method. \emph{The Annals of Mathematical Statistics} 22 (1951), 400--407.
\bibitem{robbins1956} H. Robbins. An empirical Bayes approach to statistics. In \emph{Proceedings of the Third Berkeley Symposium on Mathematical Statistics and Probability}, vol. 1, pages 157--163. University of California Press, 1956.
\bibitem{robertsrosenthal2004} G. O. Roberts and J. S. Rosenthal. General state space Markov chains and MCMC algorithms. \emph{Probability Surveys} 1 (2004), 20--71.
\bibitem{robinson2016} J. C. Robinson, J. L. Rodrigo and W. Sadowski. \emph{The three dimensional Navier Stokes equations: classical theory}. Cambridge University Press, 2016.
\bibitem{rodgerstao2020} B. Rodgers and T. Tao. The de Bruijn Newman constant is nonnegative. \emph{Forum of Mathematics, Pi} 8 (2020), e6.
\bibitem{rosenblatt1952} M. Rosenblatt. Remarks on a multivariate transformation. \emph{The Annals of Mathematical Statistics} 23 (1952), 470--472.
\bibitem{rosenthal1913} A. Rosenthal. Beweis der Unm\"oglichkeit ergodischer Gassysteme. \emph{Annalen der Physik} 42 (1913), 796--806.
\bibitem{rosinski2001} J. Rosi\'nski. Series representations of L\'evy processes from the perspective of point processes. In \emph{L\'evy processes: theory and applications}, pages 401--415. Birkh\"auser, 2001.
\bibitem{rubin1981} D. B. Rubin. The Bayesian bootstrap. \emph{The Annals of Statistics} 9 (1981), 130--134.
\bibitem{sagawaueda2010} T. Sagawa and M. Ueda. Generalized Jarzynski equality under nonequilibrium feedback control. \emph{Physical Review Letters} 104 (2010), 090602.
\bibitem{saintraymond2003} L. Saint Raymond. Convergence of solutions to the Boltzmann equation in the incompressible Euler limit. \emph{Archive for Rational Mechanics and Analysis} 166 (2003), 47--80.
\bibitem{saintraymond2009} L. Saint Raymond. \emph{Hydrodynamic limits of the Boltzmann equation}. Lecture Notes in Mathematics 1971, Springer, 2009.
\bibitem{sanov1957} I. N. Sanov. On the probability of large deviations of random variables. \emph{Matematicheskii Sbornik} 42 (1957), 11--44.
\bibitem{sarkka2011} S. S\"arkk\"a. Linear operators and stochastic partial differential equations in Gaussian process regression. In \emph{Artificial neural networks and machine learning: ICANN 2011}, Lecture Notes in Computer Science 6792, pages 151--158. Springer, 2011.
\bibitem{sato1999} K. Sato. \emph{L\'evy Processes and Infinitely Divisible Distributions}. Cambridge Studies in Advanced Mathematics 68, Cambridge University Press, 1999.
\bibitem{sauermajer2009} T. Sauer and U. Majer, eds. \emph{David Hilbert's Lectures on the Foundations of Physics 1915 to 1927: Relativity, Quantum Theory and Epistemology}. Springer, 2009.
\bibitem{savage1954} L. J. Savage. \emph{The foundations of statistics}. Wiley, 1954.
\bibitem{schilder1966} M. Schilder. Some asymptotic formulas for Wiener integrals. \emph{Transactions of the American Mathematical Society} 125 (1966), 63--85.
\bibitem{schilling2012} R. L. Schilling, R. Song and Z. Vondra\v{c}ek. \emph{Bernstein Functions: Theory and Applications}. Second edition, De Gruyter Studies in Mathematics 37, De Gruyter, 2012.
\bibitem{schmidthieber2021} J. Schmidt Hieber. The Kolmogorov Arnold representation theorem revisited. \emph{Neural Networks} 137 (2021), 119--126.
\bibitem{schnorr1973} C. P. Schnorr. Process complexity and effective random tests. \emph{Journal of Computer and System Sciences} 7 (1973), 376--388.
\bibitem{schottky1918} W. Schottky. \"Uber spontane Stromschwankungen in verschiedenen Elektrizit\"atsleitern. \emph{Annalen der Physik} 57 (1918), 541--567.
\bibitem{schrodinger1931} E. Schr\"odinger. \"Uber die Umkehrung der Naturgesetze. \emph{Sitzungsberichte der Preussischen Akademie der Wissenschaften, Physikalisch Mathematische Klasse} (1931), 144--153.
\bibitem{schwarz1978} G. Schwarz. Estimating the dimension of a model. \emph{The Annals of Statistics} 6 (1978), 461--464.
\bibitem{sedkipudlo2012} M. Sedki and P. Pudlo. Contribution to the discussion of Fearnhead and Prangle (2012): summary statistics selection by least angle regression. \emph{Journal of the Royal Statistical Society B} 74 (2012).
\bibitem{serfaty2020} S. Serfaty. Mean field limit for Coulomb type flows. \emph{Duke Mathematical Journal} 169 (2020), 2887--2935.
\bibitem{sethuraman1994} J. Sethuraman. A constructive definition of Dirichlet priors. \emph{Statistica Sinica} 4 (1994), 639--650.
\bibitem{shafervovk2006} G. Shafer and V. Vovk. The sources of Kolmogorov's Grundbegriffe. \emph{Statistical Science} 21 (2006), 70--98.
\bibitem{shakhov1968} E. M. Shakhov. Generalization of the Krook kinetic relaxation equation. \emph{Fluid Dynamics} 3 (1968), 95--96.
\bibitem{shalm2015} L. K. Shalm et al. Strong loophole free test of local realism. \emph{Physical Review Letters} 115 (2015), 250402.
\bibitem{shannon1948} C. E. Shannon. A mathematical theory of communication. \emph{Bell System Technical Journal} 27 (1948), 379--423, 623--656.
\bibitem{sheleveque1994} Z. S. She and E. L\'ev\^eque. Universal scaling laws in fully developed turbulence. \emph{Physical Review Letters} 72 (1994), 336--339.
\bibitem{silver2018} D. Silver, T. Hubert, J. Schrittwieser, I. Antonoglou, M. Lai, A. Guez, M. Lanctot, L. Sifre, D. Kumaran, T. Graepel, T. Lillicrap, K. Simonyan and D. Hassabis. A general reinforcement learning algorithm that masters chess, shogi, and Go through self play. \emph{Science} 362 (2018), 1140--1144.
\bibitem{simanyi2013} N. Sim\'anyi. Singularities and nonhyperbolic manifolds do not coincide. \emph{Nonlinearity} 26 (2013), 1703--1717.
\bibitem{simanyi2003} N. Sim\'anyi. Proof of the Boltzmann Sinai ergodic hypothesis for typical hard disk systems. \emph{Inventiones Mathematicae} 154 (2003), 123--178.
\bibitem{simanyi2004} N. Sim\'anyi. Proof of the ergodic hypothesis for typical hard ball systems. \emph{Annales Henri Poincar\'e} 5 (2004), 203--233.
\bibitem{simanyi2009} N. Sim\'anyi. Conditional proof of the Boltzmann Sinai ergodic hypothesis. \emph{Inventiones Mathematicae} 177 (2009), 381--413.
\bibitem{sinkhorn1964} R. Sinkhorn. A relationship between arbitrary positive matrices and doubly stochastic matrices. \emph{The Annals of Mathematical Statistics} 35 (1964), 876--879.
\bibitem{skitovich1953} V. P. Skitovich. On a property of the normal distribution. \emph{Doklady Akademii Nauk SSSR} 89 (1953), 217--219.
\bibitem{skorokhod1956} A. V. Skorokhod. Limit theorems for stochastic processes. \emph{Theory of Probability and Its Applications} 1 (1956), 261--290.
\bibitem{solomonoff1964} R. J. Solomonoff. A formal theory of inductive inference. Parts I and II. \emph{Information and Control} 7 (1964), 1--22, 224--254.
\bibitem{song2021} Y. Song, J. Sohl Dickstein, D. P. Kingma, A. Kumar, S. Ermon and B. Poole. Score based generative modeling through stochastic differential equations. \emph{International Conference on Learning Representations} (2021).
\bibitem{spantini2015} A. Spantini, A. Solonen, T. Cui, J. Martin, L. Tenorio and Y. Marzouk. Optimal low rank approximations of Bayesian linear inverse problems. \emph{SIAM Journal on Scientific Computing} 37 (2015), A2451--A2487.
\bibitem{spantini2018} A. Spantini, D. Bigoni and Y. Marzouk. Inference via low dimensional couplings. \emph{Journal of Machine Learning Research} 19 (2018), 1--71.
\bibitem{spohn1978} H. Spohn. The Lorentz process converges to a random flight process. \emph{Communications in Mathematical Physics} 60 (1978), 277--290.
\bibitem{spohn1991} H. Spohn. \emph{Large scale dynamics of interacting particles}. Springer, 1991.
\bibitem{stam1959} A. J. Stam. Some inequalities satisfied by the quantities of information of Fisher and Shannon. \emph{Information and Control} 2 (1959), 101--112.
\bibitem{stein1965} R. B. Stein. A theoretical analysis of neuronal variability. \emph{Biophysical Journal} 5 (1965), 173--194.
\bibitem{stentoft2004} L. Stentoft. Convergence of the least squares Monte Carlo approach to American option valuation. \emph{Management Science} 50 (2004), 1193--1203.
\bibitem{stone1982} C. J. Stone. Optimal global rates of convergence for nonparametric regression. \emph{The Annals of Statistics} 10 (1982), 1040--1053.
\bibitem{storvik2002} G. Storvik. Particle filters for state space models with the presence of unknown static parameters. \emph{IEEE Transactions on Signal Processing} 50 (2002), 281--289.
\bibitem{streater1964} R. F. Streater and A. S. Wightman. \emph{PCT, Spin and Statistics, and All That}. W. A. Benjamin, 1964.
\bibitem{stuart2010} A. M. Stuart. Inverse problems: a Bayesian perspective. \emph{Acta Numerica} 19 (2010), 451--559.
\bibitem{stuartteckentrup2018} A. M. Stuart and A. L. Teckentrup. Posterior consistency for Gaussian process approximations of Bayesian posterior distributions. \emph{Mathematics of Computation} 87 (2018), 721--753.
\bibitem{sturm2006} K. T. Sturm. On the geometry of metric measure spaces. I. \emph{Acta Mathematica} 196 (2006), 65--131.
\bibitem{swendsen1987} R. H. Swendsen and J. S. Wang. Nonuniversal critical dynamics in Monte Carlo simulations. \emph{Physical Review Letters} 58 (1987), 86--88.
\bibitem{szilard1929} L. Szilard. \"Uber die Entropieverminderung in einem thermodynamischen System bei Eingriffen intelligenter Wesen. \emph{Zeitschrift f\"ur Physik} 53 (1929), 840--856.
\bibitem{sznitman1991} A. S. Sznitman. Topics in propagation of chaos. In \emph{\'Ecole d'\'Et\'e de Probabilit\'es de Saint Flour XIX 1989}, pages 165--251. Lecture Notes in Mathematics 1464, Springer, 1991.
\bibitem{talagrand1996} M. Talagrand. Transportation cost for Gaussian and other product measures. \emph{Geometric and Functional Analysis} 6 (1996), 587--600.
\bibitem{tao2009} T. Tao. Global regularity for a logarithmically supercritical hyperdissipative Navier Stokes equation. \emph{Analysis \& PDE} 2 (2009), 361--366.
\bibitem{tao2016} T. Tao. Finite time blowup for an averaged three dimensional Navier Stokes equation. \emph{Journal of the American Mathematical Society} 29 (2016), 601--674.
\bibitem{taqqu2001} M. S. Taqqu. Bachelier and his times: a conversation with Bernard Bru. \emph{Finance and Stochastics} 5 (2001), 3--32.
\bibitem{taroni1998} F. Taroni, C. Champod and P. Margot. Forerunners of Bayesianism in early forensic science. \emph{Jurimetrics Journal} 38 (1998), 183--200.
\bibitem{tavare1997} S. Tavar\'e, D. J. Balding, R. C. Griffiths and P. Donnelly. Inferring coalescence times from DNA sequence data. \emph{Genetics} 145 (1997), 505--518.
\bibitem{taylorgreen1937} G. I. Taylor and A. E. Green. Mechanism of the production of small eddies from large ones. \emph{Proceedings of the Royal Society of London A} 158 (1937), 499--521.
\bibitem{teicher1961} H. Teicher. Maximum likelihood characterization of distributions. \emph{The Annals of Mathematical Statistics} 32 (1961), 1214--1222.
\bibitem{thorin1977} O. Thorin. On the infinite divisibility of the lognormal distribution. \emph{Scandinavian Actuarial Journal} (1977), 121--148.
\bibitem{thorin1977b} O. Thorin. On the infinite divisibility of the Pareto distribution. \emph{Scandinavian Actuarial Journal} (1977), 31--40.
\bibitem{thorin1978} O. Thorin. An extension of the notion of a generalized $\Gamma$ convolution. \emph{Scandinavian Actuarial Journal} (1978), 141--149.
\bibitem{tropp2015} J. A. Tropp. An introduction to matrix concentration inequalities. \emph{Foundations and Trends in Machine Learning} 8 (2015), 1--230.
\bibitem{tsai1998} T. P. Tsai. On Leray's self similar solutions of the Navier Stokes equations satisfying local energy estimates. \emph{Archive for Rational Mechanics and Analysis} 143 (1998), 29--51.
\bibitem{tsirelson1980} B. S. Tsirelson. Quantum generalizations of Bell's inequality. \emph{Letters in Mathematical Physics} 4 (1980), 93--100.
\bibitem{tsybakov2009} A. B. Tsybakov. \emph{Introduction to nonparametric estimation}. Springer, 2009.
\bibitem{tuowu2016} R. Tuo and C. F. J. Wu. A theoretical framework for calibration in computer models: parametrization, estimation and convergence properties. \emph{SIAM/ASA Journal on Uncertainty Quantification} 4 (2016), 767--795.
\bibitem{uhlmann1977} A. Uhlmann. Relative entropy and the Wigner Yanase Dyson Lieb concavity in an interpolation theory. \emph{Communications in Mathematical Physics} 54 (1977), 21--32.
\bibitem{ulamvonneumann1947} S. M. Ulam and J. von Neumann. On combination of stochastic and deterministic processes (abstract). \emph{Bulletin of the American Mathematical Society} 53 (1947), 1120.
\bibitem{vakeroudis2012} S. Vakeroudis. Bougerol's identity in law and extensions. \emph{Probability Surveys} 9 (2012), 411--437.
\bibitem{vancampenhoutcover1981} J. Van Campenhout and T. M. Cover. Maximum entropy and conditional probability. \emph{IEEE Transactions on Information Theory} 27 (1981), 483--489.
\bibitem{vandervaart1998} A. W. van der Vaart. \emph{Asymptotic statistics}. Cambridge University Press, 1998.
\bibitem{vandervaartwellner2023} A. W. van der Vaart and J. A. Wellner. \emph{Weak convergence and empirical processes: with applications to statistics}. Second edition, Springer, 2023.
\bibitem{varadarajan1985} V. S. Varadarajan. \emph{Geometry of Quantum Theory}. Second edition, Springer, 1985.
\bibitem{vasseur2007} A. F. Vasseur. A new proof of partial regularity of solutions to Navier Stokes equations. \emph{Nonlinear Differential Equations and Applications} 14 (2007), 753--785.
\bibitem{vehtari2021} A. Vehtari, A. Gelman, D. Simpson, B. Carpenter and P. C. B\"urkner. Rank normalization, folding, and localization: an improved $\widehat R$ for assessing convergence of MCMC (with discussion). \emph{Bayesian Analysis} 16 (2021), 667--718.
\bibitem{vershynin2018} R. Vershynin. \emph{High dimensional probability}. Cambridge University Press, 2018.
\bibitem{villani2002review} C. Villani. A review of mathematical topics in collisional kinetic theory. In \emph{Handbook of mathematical fluid dynamics}, vol. 1, pages 71--305. North Holland, 2002.
\bibitem{villani2003} C. Villani. \emph{Topics in optimal transportation}. American Mathematical Society, 2003.
\bibitem{villani2003cercignani} C. Villani. Cercignani's conjecture is sometimes true and always almost true. \emph{Communications in Mathematical Physics} 234 (2003), 455--490.
\bibitem{villani2009} C. Villani. \emph{Optimal transport: old and new}. Springer, 2009.
\bibitem{vonneumann1927a} J. von Neumann. Wahrscheinlichkeitstheoretischer Aufbau der Quantenmechanik. \emph{Nachrichten von der Gesellschaft der Wissenschaften zu G\"ottingen, Mathematisch Physikalische Klasse} (1927), 245--272.
\bibitem{vonneumann1927b} J. von Neumann. Thermodynamik quantenmechanischer Gesamtheiten. \emph{Nachrichten von der Gesellschaft der Wissenschaften zu G\"ottingen, Mathematisch Physikalische Klasse} (1927), 273--291.
\bibitem{vonneumann1930} J. von Neumann. Allgemeine Eigenwerttheorie Hermitescher Funktionaloperatoren. \emph{Mathematische Annalen} 102 (1930), 49--131.
\bibitem{vonneumann1932} J. von Neumann. \emph{Mathematische Grundlagen der Quantenmechanik}. Springer, 1932.
\bibitem{vonneumann1932ergodic} J. von Neumann. Proof of the quasiergodic hypothesis. \emph{Proceedings of the National Academy of Sciences} 18 (1932), 70--82.
\bibitem{vonneumann1932operator} J. von Neumann. Zur Operatorenmethode in der klassischen Mechanik. \emph{Annals of Mathematics} 33 (1932), 587--642.
\bibitem{vonplato1994} J. von Plato. \emph{Creating modern probability: its mathematics, physics and philosophy in historical perspective}. Cambridge University Press, 1994.
\bibitem{vonrenesse2005} M. K. von Renesse and K. T. Sturm. Transport inequalities, gradient estimates, entropy and Ricci curvature. \emph{Communications on Pure and Applied Mathematics} 58 (2005), 923--940.
\bibitem{smoluchowski1906} M. von Smoluchowski. Zur kinetischen Theorie der Brownschen Molekularbewegung und der Suspensionen. \emph{Annalen der Physik} 21 (1906), 756--780.
\bibitem{smoluchowski1912} M. von Smoluchowski. Experimentell nachweisbare, der \"ublichen Thermodynamik widersprechende Molekularph\"anomene. \emph{Physikalische Zeitschrift} 13 (1912), 1069--1080.
\bibitem{wallstrom1994} T. C. Wallstrom. Inequivalence between the Schr\"odinger equation and the Madelung hydrodynamic equations. \emph{Physical Review A} 49 (1994), 1613--1617.
\bibitem{walters1982} P. Walters. \emph{An introduction to ergodic theory}. Graduate Texts in Mathematics 79, Springer, 1982.
\bibitem{wangrockova2020} Y. Wang and V. Ro\v{c}kov\'a. Uncertainty quantification for sparse deep learning. In \emph{Proceedings of the Twenty Third International Conference on Artificial Intelligence and Statistics}, PMLR 108, pages 298--308, 2020.
\bibitem{wegmann2009} D. Wegmann, C. Leuenberger and L. Excoffier. Efficient approximate Bayesian computation coupled with Markov chain Monte Carlo without likelihood. \emph{Genetics} 182 (2009), 1207--1218.
\bibitem{west1981} M. West. Robust sequential approximate Bayesian estimation. \emph{Journal of the Royal Statistical Society B} 43 (1981), 157--166.
\bibitem{west1993} M. West. Approximating posterior distributions by mixtures. \emph{Journal of the Royal Statistical Society B} 55 (1993), 409--422.
\bibitem{weyl1916} H. Weyl. \"Uber die Gleichverteilung von Zahlen mod. Eins. \emph{Mathematische Annalen} 77 (1916), 313--352.
\bibitem{whittle1954} P. Whittle. On stationary processes in the plane. \emph{Biometrika} 41 (1954), 434--449.
\bibitem{whittle1955} P. Whittle. Some distribution and moment formulae for the Markov chain. \emph{Journal of the Royal Statistical Society B} 17 (1955), 235--242.
\bibitem{whittle1963} P. Whittle. Stochastic processes in several dimensions. \emph{Bulletin of the International Statistical Institute} 40 (1963), 974--994.
\bibitem{wiener1923} N. Wiener. Differential space. \emph{Journal of Mathematics and Physics} 2 (1923), 131--174.
\bibitem{wiener1948} N. Wiener. \emph{Cybernetics or control and communication in the animal and the machine}. The Technology Press and John Wiley \& Sons, 1948.
\bibitem{williams2006} B. Williams, D. Higdon, J. Gattiker, L. Moore, M. D. McKay and S. Keller McNulty. Combining experimental data and computer simulations, with an application to flyer plate experiments. \emph{Bayesian Analysis} 1 (2006), 765--792.
\bibitem{wonham1964} W. M. Wonham. Some applications of stochastic differential equations to optimal nonlinear filtering. \emph{Journal of the Society for Industrial and Applied Mathematics, Series A: Control} 2 (1964), 347--369.
\bibitem{xiebarron2000} Q. Xie and A. R. Barron. Asymptotic minimax regret for data compression, gambling, and prediction. \emph{IEEE Transactions on Information Theory} 46 (2000), 431--445.
\bibitem{xiong2008} J. Xiong. \emph{An introduction to stochastic filtering theory}. Oxford University Press, 2008.
\bibitem{khinchin1938} A. Ya. Khinchin. On unimodal distributions. \emph{Izvestiya Nauchno Issledovatelskogo Instituta Matematiki i Mekhaniki Tomskogo Gosudarstvennogo Universiteta} 2 (1938), 1--7. In Russian.
\bibitem{sinai1959} Ya. G. Sinai. On the notion of entropy of a dynamical system. \emph{Doklady Akademii Nauk SSSR} 124 (1959), 768--771.
\bibitem{sinai1970} Ya. G. Sinai. Dynamical systems with elastic reflections: ergodic properties of dispersing billiards. \emph{Russian Mathematical Surveys} 25 (1970), 137--189.
\bibitem{pesin1977} Ya. B. Pesin. Characteristic Lyapunov exponents and smooth ergodic theory. \emph{Russian Mathematical Surveys} 32 (1977), 55--114.
\bibitem{sinai1992} Ya. G. Sinai. Statistics of shocks in solutions of inviscid Burgers equation. \emph{Communications in Mathematical Physics} 148 (1992), 601--621.
\bibitem{yamazato1978} M. Yamazato. Unimodality of infinitely divisible distribution functions of class L. \emph{The Annals of Probability} 6 (1978), 523--531.
\bibitem{yang1952} C. N. Yang. The spontaneous magnetization of a two dimensional Ising model. \emph{Physical Review} 85 (1952), 808--816.
\bibitem{yau1991} H. T. Yau. Relative entropy and hydrodynamics of Ginzburg Landau models. \emph{Letters in Mathematical Physics} 22 (1991), 63--80.
\bibitem{yedidia2005} J. S. Yedidia, W. T. Freeman and Y. Weiss. Constructing free energy approximations and generalized belief propagation algorithms. \emph{IEEE Transactions on Information Theory} 51 (2005), 2282--2312.
\bibitem{prokhorov1956} Yu. V. Prokhorov. Convergence of random processes and limit theorems in probability theory. \emph{Theory of Probability and Its Applications} 1 (1956), 157--214.
\bibitem{shtarkov1987} Yu. M. Shtarkov. Universal sequential coding of single messages. \emph{Problems of Information Transmission} 23 (1987), 175--186.
\bibitem{yuzvinskii1967} S. A. Yuzvinskii. Computing the entropy of a group of endomorphisms. \emph{Siberian Mathematical Journal} 8 (1967), 172--178.
\bibitem{zakai1969} M. Zakai. On the optimal filtering of diffusion processes. \emph{Zeitschrift f\"ur Wahrscheinlichkeitstheorie und verwandte Gebiete} 11 (1969), 230--243.
\bibitem{zdeborova2016} L. Zdeborov\'a and F. Krzakala. Statistical physics of inference: thresholds and algorithms. \emph{Advances in Physics} 65 (2016), 453--552.
\bibitem{zermelo1896} E. Zermelo. Ueber einen Satz der Dynamik und die mechanische W\"armetheorie. \emph{Annalen der Physik} 57 (1896), 485--494.
\bibitem{zvonkinlevin1970} A. K. Zvonkin and L. A. Levin. The complexity of finite objects and the development of the concepts of information and randomness by means of the theory of algorithms. \emph{Russian Mathematical Surveys} 25 (1970), 83--124.

\end{thebibliography}
